% formalisation.tex: blueprint extension of an accepted paper (AGENTS.md, section 6.6, case c).
% Not refereed: Lean checks every node. paper.tex is never edited.
% Nodes of paper.tex are referred to by their labels: xr-hyper reads paper.aux, so compile
% paper.tex first (CI does). New labels start with f (flem:, fprop:, fdef:) and are never reused.
% Every node has \lean{...} with its Lean name and an exhaustive \uses{...}; never write \leanok.
% Compile: latexmk -pdf -interaction=nonstopmode -halt-on-error -cd papers/<P>/formalisation.tex
\documentclass[11pt]{article}
\usepackage[T1]{fontenc}
\usepackage[utf8]{inputenc}
\usepackage{lmodern}
\usepackage[a4paper,margin=2.5cm]{geometry}
\usepackage{amsmath,amssymb,mathtools}
\usepackage{xcolor}
\usepackage{xr-hyper}
\usepackage[colorlinks=true,linkcolor=blue!50!black,citecolor=blue!50!black,urlcolor=blue!50!black]{hyperref}
\usepackage{../../templates/paper/blueprint-macros}
\externaldocument{paper}
\counterwithout{theorem}{section}
\renewcommand{\thetheorem}{F\arabic{theorem}}

% Notation of the paper, written out.
\newcommand{\E}{\mathbb{E}}
\newcommand{\PP}{\mathbb{P}}
\newcommand{\KL}{\mathrm{KL}}
\newcommand{\Kinf}{\mathcal K_{\inf}}
\newcommand{\cP}{\mathcal P}

\title{Formalisation notes for P000-hy-unknown-imposed-grid}
\author{Incognita agents\thanks{Not refereed: every node is checked in Lean. Numbers F1, F2, \dots{}
  are this file's; other numbers refer to the paper.}}
\date{2026-11-17}

\begin{document}
\maketitle

Rounds are numbered from $0$ (the paper's round $t$ is the Lean round $t-1$); the count $N_c(t)$
is the number of pulls of $c$ in the rounds $0,\dots,t-1$, so the quantities at the batch ends
$t_\ell$ do not shift. A policy ``sees the rewards before $r$'' at the round $n$ when it takes the
same value on two histories of the rounds $<n$ with the same actions and the same rewards at the
rounds $<r$; a learner batched on a grid sees, at a round of batch $\ell+1$ (the rounds
$t_\ell\le n<t_{\ell+1}$), the rewards before $t_\ell$. $\Kinf(\nu_a,\mu)$ is the primal formula
$\inf\{\KL(\nu_a,Q):\ Q\text{ a law on }[0,1],\ \E_Q[X]>\mu\}$, as in the frozen statements.

\section{Extra lemmas}

\subsection{One batch ahead}

\begin{lemma}[masked histories]\label{flem:masked-law}
\lean{Incognita.BatchedBandits.map_history_ite_lt_eq, Learning.IsAlgEnvSeq.isAlgEnvSeqUntil_ite_lt}
\uses{def:learner}
Let $r\le M$, and let a learner see, at every round $n<M$, the rewards before $r$. Let two
environments have the same feedback kernels at the rounds $<r$. Then the history of the rounds
$<M$ in which the rewards of the rounds $\ge r$ are replaced by a fixed value $y_0$ has the same
law under the two environments. In particular the actions of the rounds $<M$, and every count
$N_c(M)$, have the same law.
\end{lemma}

\begin{proof}
\uses{def:learner}
The masked process (observations, actions, and rewards replaced by $y_0$ at the rounds $\ge r$) is
a run until $M$ of the learner against the environment $e_M$ whose feedback kernel is the one of
the given environment at the rounds $<r$ and $\delta_{y_0}$ at the rounds $\ge r$: the conditional
law of the action at a round $n<M$ given the masked history is the policy, since the policy only
reads the actions and the rewards before $r$, which the mask keeps; at the rounds $<r$ nothing is
masked; at the rounds $\ge r$ the masked reward is constant. The same holds for the second
environment with the same $e_M$, and the law of the history of a run until $M$ only depends on
the learner and on the environment.
\end{proof}

\begin{lemma}[divergence for the switched environment]\label{flem:switch-kl}
\lean{Incognita.BatchedBandits.klDiv_map_hist_switch}
\uses{def:learner}
Let $e$ be the environment that follows the bandit $\nu'$ at the rounds $<r$ and the bandit $\nu$
at the rounds $\ge r$, and $r\le M$. The relative entropy between the laws of the histories of the
rounds $<M$ under $\nu$ and under $e$ equals the relative entropy between the laws of the histories
of the rounds $<r$ under $\nu$ and under $\nu'$.
\end{lemma}

\begin{proof}
\uses{def:learner}
Chain rule over the rounds (Lattimore and Szepesv\'ari, \emph{Bandit Algorithms}, Lemma~15.1): the rounds $\ge r$ have the same step kernel
under $\nu$ and under $e$ and contribute $0$; the rounds $<r$ have the step kernels of $\nu'$.
\end{proof}

\begin{lemma}[finite horizon, one batch ahead]\label{flem:onebatch-finite}
\lean{Incognita.BatchedBandits.exp_neg_le_of_seesRewardsUpTo}
\uses{def:learner}
Let $r\le M$, $M>0$, let a learner see the rewards before $r$ at every round $n<M$, and let the
bandits $\nu,\nu'$ differ only at the arm $a$, with $\KL(\nu_a,\nu'_a)<\infty$. Then
\[
\tfrac12\exp\big(-\E_\nu[N_a(r)]\,\KL(\nu_a,\nu'_a)\big)
\ \le\ \frac{2\,\E_\nu[N_a(M)]}{M}+\frac{2\,(M-\E_{\nu'}[N_a(M)])}{M}.
\]
\end{lemma}

\begin{proof}
\uses{flem:masked-law, flem:switch-kl, def:learner}
Let $e$ be the switched environment of Lemma~\ref{flem:switch-kl}. By
Lemma~\ref{flem:masked-law} (the environments $e$ and $\nu'$ agree before $r$), the masked history
of the rounds $<M$ has the same law $\mu'$ under $e$ and under $\nu'$; let $\mu$ be its law under
$\nu$. By data processing and Lemma~\ref{flem:switch-kl},
$\KL(\mu,\mu')\le\KL(\nu\text{-history before }r,\ \nu'\text{-history before }r)
=\E_\nu[N_a(r)]\KL(\nu_a,\nu'_a)$ (divergence decomposition). The Bretagnolle--Huber inequality
(in Lean, LMLExplore's \texttt{LMLExplore.exp\_neg\_klDiv\_div\_two\_le}) on the event $\{2N_a(M)\ge M\}$ (a function of the masked history) gives
$\tfrac12e^{-\KL(\mu,\mu')}\le\PP_\nu(2N_a(M)\ge M)+\PP_{\nu'}(2N_a(M)<M)$; Markov's inequality,
for $N_a(M)$ under $\nu$ and for $M-N_a(M)\ge0$ under $\nu'$, concludes.
\end{proof}

\begin{lemma}[alternative laws have mean below $1$]\label{flem:mean-lt-one}
\lean{Incognita.BatchedBandits.integral_id_lt_one_of_klDiv_ne_top}
\uses{def:model}
If $Q$ is a law on $[0,1]$ and $\KL(P,Q)<\infty$ for a law $P$ of mean $<1$, then $Q$ has mean
$<1$.
\end{lemma}

\begin{proof}
\uses{def:model}
A law on $[0,1]$ of mean $1$ is $\delta_1$, and $P\ll\delta_1$ would give $P=\delta_1$.
\end{proof}

\subsection{HY is batched and does not know the grid ahead}

\begin{lemma}[HY sees the rewards of the previous batches]\label{flem:hy-congr}
\lean{Incognita.BatchedBandits.HY.next_congr, Incognita.BatchedBandits.isBatched_hyAlg}
\uses{def:hy}
At a round of batch $\ell+1$, two histories with the same actions and the same rewards at the
rounds $<t_\ell$ give the same action of HY$(w_0,\beta,\omega,\vartheta,\gamma)$, for all
parameters. Hence HY is batched on its grid.
\end{lemma}

\begin{proof}
\uses{def:hy}
Every statistic of Definition~\ref{def:hy} used in batch $\ell+1$ is computed at the cutoffs
$t_\ell$ and $t_{\ell-1}\le t_\ell$ (counts, pseudo-means, divergences of the pseudo-empirical laws,
leader, reference, threshold, targets), or is a current count $N_c(t-1)$, which only uses the
actions. Each is checked to agree on the two histories, in the order of the definition.
\end{proof}

\begin{lemma}[HY does not know the grid ahead]\label{flem:hy-grid}
\lean{Incognita.BatchedBandits.HY.next_hy_eq_of_t_eq}
\uses{def:hy, def:grid}
With the default parameters $\omega^{\mathrm{dy}},\vartheta,\gamma$, at a round $n$ of batch
$\ell+1$ of two grids with the same ends $t_0,\dots,t_\ell$, HY takes the same action.
\end{lemma}

\begin{proof}
\uses{def:hy, def:grid}
The grid enters the action of batch $\ell+1$ only through $t_\ell$, $t_{\ell-1}$, $x_\ell$,
$x_{\ell-1}$, $\mathsf L_\ell$, $\mathsf L_1$ (for $\ell=1$), $U_{\ell-1},U_\ell,U_{\ell+1}$
($U_{j+1}$ is a function of $\mathsf L_j$), $s_\ell$, $\omega^{\mathrm{dy}}_\ell$, $\vartheta_\ell$ and
$\gamma_\ell$, all functions of $t_0,\dots,t_\ell$; batch $1$ does not use the grid.
\end{proof}

\subsection{Assembling the lower bounds}

\begin{lemma}[$\liminf$ of a weighted sum]\label{flem:liminf-sum}
\lean{Incognita.BatchedBandits.eventually_lt_sum_mul}
\uses{}
For a finite set $S$, weights $w_i>0$ and sequences with $\liminf_\ell u_i(\ell)\ge L_i$,
$\liminf_\ell\sum_{i\in S}w_iu_i(\ell)\ge\sum_{i\in S}w_iL_i$.
\end{lemma}

\begin{proof}
\uses{}
Induction on $S$, splitting the margin in two halves.
\end{proof}

\begin{lemma}[$\liminf$ of a product]\label{flem:liminf-mul}
\lean{Incognita.BatchedBandits.eventually_lt_mul_of_tendsto}
\uses{}
If $f\ge0$ eventually, $\liminf f\ge A$ and $g\to B>0$, then $\liminf fg\ge AB$.
\end{lemma}

\begin{proof}
\uses{}
For $0\le c<AB$, choose $A'<A$ and $B'<B$ positive with $c<A'B'$.
\end{proof}

\begin{lemma}[two ratios]\label{flem:log-ratio}
\lean{Incognita.BatchedBandits.Grid.tendsto_log_t_add_two_div_log_t}
\uses{def:grid, def:ass-x}
On a grid with $x_\ell\to x$, $\log t_{\ell+2}/\log t_\ell=x_{\ell+1}x_{\ell+2}\to x^2$.
\end{lemma}

\begin{proof}
\uses{def:grid, def:ass-x}
For $\ell\ge1$ the identity holds by definition of the ratios.
\end{proof}

\begin{lemma}[Case A of Theorem~\ref{thm:lower}(ii)]\label{flem:caseA}
\lean{Incognita.BatchedBandits.eventually_lt_div_log_t_of_eventually_lt_div_log_t_add_two}
\uses{def:grid, def:ass-x}
On a grid with $x_\ell\to x>0$, a nonnegative sequence with $\liminf_\ell u_\ell/\log t_{\ell+2}\ge A$
has $\liminf_\ell u_\ell/\log t_\ell\ge x^2A$.
\end{lemma}

\begin{proof}
\uses{flem:liminf-mul, flem:log-ratio, def:grid, def:ass-x}
$u_\ell/\log t_\ell=(u_\ell/\log t_{\ell+2})\cdot(\log t_{\ell+2}/\log t_\ell)$.
\end{proof}

\begin{lemma}[Case B of Theorem~\ref{thm:lower}(ii)]\label{flem:caseB}
\lean{Incognita.BatchedBandits.frequently_lt_div_log_t_of_not_isBigO}
\uses{def:grid}
A nonnegative sequence that is not $O(\log t_\ell)$ has $\limsup_\ell u_\ell/\log t_\ell=\infty$.
\end{lemma}

\begin{proof}
\uses{def:grid}
Otherwise $u_\ell\le c\log t_\ell$ eventually, and $\log t_\ell>0$ for $\ell\ge1$.
\end{proof}

\begin{lemma}[the regret with the constant $C_1$, on every grid]\label{flem:regret-one}
\lean{Incognita.BatchedBandits.eventually_lt_regret_div_log_t}
\uses{def:grid, def:model, def:learner, def:kinf}
For every grid, every batched learner consistent on $\cP$ and every $\nu\in\cP$,
$\liminf_\ell\mathfrak R(t_\ell)/\log t_\ell\ge\sum_{c:\Delta_c>0}\Delta_c/\Kinf(\nu_c,\mathbf m^\star)$.
\end{lemma}

\begin{proof}
\uses{flem:liminf-sum, lem:onebatch, def:grid, def:model, def:learner, def:kinf}
Lemma~\ref{flem:liminf-sum} with $w_c=\Delta_c$ and, for each sub-optimal arm $c$, the bound
$\liminf_\ell\E_\nu[N_c(t_\ell)]/\log t_\ell\ge1/\Kinf(\nu_c,\mathbf m^\star)$ on every grid, which
follows from Lemma~\ref{lem:onebatch} and $\log t_\ell\le\log t_{\ell+1}$ (the argument of the
formal route to Theorem~\ref{thm:lower}(i) below); the optimal arms have $\Delta_c=0$.
\end{proof}

\begin{lemma}[the regret when some count is not logarithmic]\label{flem:regret-caseB}
\lean{Incognita.BatchedBandits.frequently_lt_regret_div_log_t_of_not_isBigO}
\uses{def:grid, def:learner}
If a sub-optimal arm $c$ has $\E_\nu[N_c(t_\ell)]\ne O(\log t_\ell)$, then
$\limsup_\ell\mathfrak R(t_\ell)/\log t_\ell=\infty$.
\end{lemma}

\begin{proof}
\uses{flem:caseB, def:grid, def:learner}
$\mathfrak R(t_\ell)\ge\Delta_c\E_\nu[N_c(t_\ell)]$, and Lemma~\ref{flem:caseB}.
\end{proof}

\subsection{Two batches ahead, without concentration}

The paper proves Lemma~\ref{lem:twobatches} with boundary tilts, a tilted measure on the reward
stacks and the typical values of the log-likelihood ratios (Lemmas~\ref{lem:tilts}--\ref{lem:typ}).
The Lean proof of part (ii) uses instead a one-sided change of measure between two alternatives
$\kappa^\pm$ with $\kappa^-_a\ge(1-s)\kappa^+_a$, and data processing from the environment that
is $\nu$ before $t_\ell$; no concentration is needed.

\begin{lemma}[one-sided change of measure]\label{flem:com-one-sided}
\lean{Incognita.BatchedBandits.withDensity_map_hist_le, Incognita.BatchedBandits.setLIntegral_stepKernel_le}
\uses{def:learner}
Let two environments $e_1,e_2$ satisfy $w_n(b)\,e_2^{(n)}(\cdot\mid h,b)\le e_1^{(n)}(\cdot\mid h,b)$
for the feedback kernels of every round $n$, history $h$ and action $b$. Then for every measurable
set $A$ of histories of the first $n$ rounds,
$\PP_{e_1}(H_n\in A)\ge\E_{e_2}\big[\mathbf 1_A(H_n)\prod_{i<n}w_i(A_i)\big]$.
\end{lemma}

\begin{proof}
\uses{def:learner}
Induction on $n$: the law of $H_{n+1}$ is the law of $H_n$ composed with the step kernel of the
round $n$, which under $e_1$ dominates the step kernel under $e_2$ with the density $w_n(A_n)$
(same observation kernel and policy).
\end{proof}

\begin{lemma}[switched environment, bandit on the right]\label{flem:switch-kl-left}
\lean{Incognita.BatchedBandits.klDiv_map_hist_switch_left, Learning.IsAlgEnvSeq.isAlgEnvSeqUntil_of_feedback_eq}
\uses{def:learner}
Let $e$ follow the bandit $\nu$ at the rounds $<r$ and $\nu'$ afterwards, and $r\le M$. The relative
entropy between the laws of the histories of the first $M$ rounds under $e$ and under $\nu'$ is the
relative entropy between the laws of the histories of the first $r$ rounds under $\nu$ and
under $\nu'$.
\end{lemma}

\begin{proof}
\uses{def:learner}
Chain rule; before $r$ the run against $e$ is a run against $\nu$, so the histories have the same
laws; after $r$ the step kernels agree.
\end{proof}

\begin{lemma}[data processing, multiplicative form]\label{flem:dp}
\lean{Incognita.BatchedBandits.exp_neg_div_le_of_klBer_le}
\uses{}
If $\mathrm{kl}(q,p)\le D$ with $q\in(0,1]$, $p\in[0,1]$, then $p\ge e^{-(D+\log2)/q}$.
\end{lemma}

\begin{proof}
\uses{}
$\mathrm{kl}(q,p)\ge q\log(1/p)-\log2$, since the binary entropy is at most $\log2$ and
$\log(1-p)\le0$.
\end{proof}

\begin{lemma}[the alternatives]\label{flem:alt-laws}
\lean{Incognita.BatchedBandits.exists_two_alternatives, Incognita.BatchedBandits.exists_unitLaws_mean_mem_Ioo, Incognita.BatchedBandits.mix_mem_unitLaws, Incognita.klDiv_mix_le}
\uses{def:model, def:kinf}
Let $\nu\in\cP$, $a$ sub-optimal, $Q$ a law on $[0,1]$ with mean $>\mathbf m^\star$ and
$\KL(\nu_a,Q)<\infty$, and $s\in(0,1)$. There are $\kappa^+,\kappa^-\in\cP$ equal to $\nu$ outside
$a$, with $a$ the unique optimal arm of $\kappa^+$, $a$ sub-optimal for $\kappa^-$,
$\kappa^-_a\ge(1-s)\kappa^+_a$ and $\KL(\nu_a,\kappa^+_a)\le\KL(\nu_a,Q)$.
\end{lemma}

\begin{proof}
\uses{flem:mean-lt-one, def:model}
$\kappa^+_a=(1-\theta)\nu_a+\theta Q$ with $\theta$ chosen so that its mean is in
$(\mathbf m^\star,\mathbf m^\star/(1-s))$; $\KL(\nu_a,\kappa^+_a)\le\theta\KL(\nu_a,Q)$ by
convexity. $\kappa^-_a=(1-s)\kappa^+_a+s\delta_0$ has mean $<\mathbf m^\star$.
\end{proof}

\begin{lemma}[finite horizon, two batches ahead]\label{flem:twobatches-finite}
\lean{Incognita.BatchedBandits.pow_mul_exp_neg_le_integral_pullCount_div_add}
\uses{def:learner}
Let $r_0\le r\le M$, $M>0$, let the learner see the rewards before $r_0$ at the rounds $<r$ and the
rewards before $r$ at the rounds $<M$, let $\nu,\kappa^+,\kappa^-$ agree outside $a$, with
$\kappa^-_a\ge c\,\kappa^+_a$, $c\in[0,1]$, $\KL(\nu_a,\kappa^+_a)<\infty$, and let
$q=\PP_\nu(N_a(r)\le L)>0$. Then
\[
c^L\exp\Big(-\frac{\E_\nu[N_a(r_0)]\KL(\nu_a,\kappa^+_a)+\log2}{q}\Big)\le
\frac{2\E_{\kappa^-}[N_a(M)]}{M}+\frac{2(M-\E_{\kappa^+}[N_a(M)])}{M}.
\]
\end{lemma}

\begin{proof}
\uses{flem:masked-law, flem:com-one-sided, flem:switch-kl-left, flem:dp, def:learner}
Let $e_Q$ follow $\nu$ before $r_0$ and $\kappa^+$ afterwards. By Lemma~\ref{flem:masked-law} (cutoff
$r_0$, horizon $r$), $\PP_{e_Q}(N_a(r)\le L)=q$; by data processing and
Lemma~\ref{flem:switch-kl-left}, $\mathrm{kl}(q,\PP_{\kappa^+}(N_a(r)\le L))\le
\E_\nu[N_a(r_0)]\KL(\nu_a,\kappa^+_a)$, so Lemma~\ref{flem:dp} bounds $\PP_{\kappa^+}(N_a(r)\le L)$
from below. Let $G=\{2N_a(M)\ge M\}$, an event of the masked history at $r$. By
Lemma~\ref{flem:masked-law} (cutoff $r$) and Lemma~\ref{flem:com-one-sided} (the environment
$\kappa^-$ before $r$, $\kappa^+$ afterwards, against $\kappa^+$, with $w_n(b)=c$ for $n<r$,
$b=a$), $\PP_{\kappa^-}(G)\ge\E_{\kappa^+}[c^{N_a(r)}\mathbf 1_G]$. Hence
$\PP_{\kappa^-}(G)+\PP_{\kappa^+}(G^c)\ge\E_{\kappa^+}[c^{N_a(r)}]\ge c^L\PP_{\kappa^+}(N_a(r)\le L)$,
and Markov's inequality concludes.
\end{proof}

\begin{lemma}[two batches ahead]\label{flem:twobatches}
\lean{Incognita.BatchedBandits.eventually_lt_integral_pullCount_div_log_t_add_two, Incognita.BatchedBandits.mul_log_t_add_two_lt_integral_pullCount, Incognita.BatchedBandits.mul_lt_of_mul_neg_add_neg_div_le, Incognita.BatchedBandits.two_mul_integral_pullCount_div_add_le_rpow_div}
\uses{def:grid, def:model, def:learner, def:kinf}
For every grid, every batched learner consistent on $\cP$, $\nu\in\cP$ and $a$ sub-optimal with
$\E_\nu[N_a(t_{\ell+1})]=O(\log t_{\ell+2})$:
$\liminf_\ell\E_\nu[N_a(t_\ell)]/\log t_{\ell+2}\ge1/\Kinf(\nu_a,\mathbf m^\star)$.
\end{lemma}

\begin{proof}
\uses{flem:twobatches-finite, flem:alt-laws, def:kinf, def:model, def:learner, def:grid}
Let $0\le c<c'<1/\Kinf$, $Q$ from the primal formula with $c'\KL(\nu_a,Q)<1$,
$\gamma=1-c'\KL(\nu_a,Q)$, $C'\ge1$ with $\E_\nu[N_a(t_{\ell+1})]\le C'\log t_{\ell+2}$ eventually,
$\lambda=\gamma^2/(16C')$, $s=1-e^{-\lambda}$, and $\kappa^\pm$ from Lemma~\ref{flem:alt-laws}. Apply
Lemma~\ref{flem:twobatches-finite} with $r_0=t_\ell$, $r=t_{\ell+1}$, $M=t_{\ell+2}$, $c=1-s$ and
$L=\lceil4\E_\nu[N_a(r)]/\gamma\rceil$, so that $q\ge1-\gamma/4$ by Markov's inequality. Consistency
under $\kappa^-$ (for $a$) and $\kappa^+$ (for the other arms) at the exponent $p=\gamma/4$ bounds
the right-hand side by $M^{p-1}$. Taking logarithms, with $L\lambda\le(\gamma/4)\log M+\lambda$,
gives $\E_\nu[N_a(t_\ell)]\KL(\nu_a,\kappa^+_a)+\log2\ge(1-3\gamma/4)\log M-\lambda$, while
$c'\KL(\nu_a,\kappa^+_a)\le1-\gamma$; hence $\E_\nu[N_a(t_\ell)]>c'\log M$ once
$(\gamma/4)\log M>\lambda+\log2+1$.
\end{proof}

\subsection{Upper bound}

\begin{lemma}[the rational and the real suprema of HY's dual formula]\label{flem:rational-sup}
\lean{Incognita.BatchedBandits.HY.pseudoKinf_eq_iSup_real, Incognita.BatchedBandits.HY.ofReal_pseudoLogMean_le_pseudoKinf, Incognita.BatchedBandits.HY.continuousAt_pseudoLogMean, Incognita.BatchedBandits.HY.pseudoKinf_le_of_forall}
\uses{def:kinf, def:hy}
Let the rewards of a history be at most $1$, $\mu\in[0,1)$, and, for an arm $c$, a cutoff $r$ and a
pseudo-count $w_0$, let $F(\lambda)=\E_{\tilde\nu_c}[\log(1-\lambda(X-\mu))]$ for the
pseudo-empirical law $\tilde\nu_c$ of the rewards of $c$ before $r$. Then
$\sup\{F(\lambda):\lambda\in[0,\frac1{1-\mu})\cap\mathbb Q\}=\sup\{F(\lambda):\lambda\in[0,\frac1{1-\mu})\}$
(both truncated at $0$): the divergence $\Kinf(\tilde\nu_c,\mu)$ that HY computes with rational
$\lambda$ is the one of Definition~\ref{def:kinf}.
\end{lemma}

\begin{proof}
\uses{def:kinf, def:hy}
On $[0,\frac1{1-\mu})$ every argument $1-\lambda(y-\mu)$ ($y\le1$) and $1+\lambda\mu$ is positive,
so $F$ is continuous there; for $\lambda$ in the interval and $a<F(\lambda)$, a rational
$\lambda'\in(\lambda,\frac1{1-\mu})$ close enough to $\lambda$ has $F(\lambda')>a$.
\end{proof}

\begin{lemma}[the statistics of HY are frozen along a run]\label{flem:frozen}
\lean{Incognita.BatchedBandits.HY.count_history, Incognita.BatchedBandits.HY.pseudoMean_history, Incognita.BatchedBandits.HY.pseudoKinf_history, Incognita.BatchedBandits.HY.leader_history, Incognita.BatchedBandits.HY.imedKinf_history, Incognita.BatchedBandits.HY.reference_history, Incognita.BatchedBandits.HY.testKinf_history, Incognita.BatchedBandits.HY.threshold_history, Incognita.BatchedBandits.HY.target_history, Incognita.BatchedBandits.HY.imedIndex_history, Incognita.BatchedBandits.HY.imedIndex_history_mono, Incognita.BatchedBandits.hyAlg_ae_action_eq, Incognita.BatchedBandits.HY.next_of_imedPhase, Incognita.BatchedBandits.HY.next_of_pacedPhase, Incognita.BatchedBandits.hyAlg_ae_imedIndex_action_le}
\uses{def:hy, def:learner, def:grid}
On a run, for $r\le n$, the statistics of HY at the cutoff $r$ computed from the history of the
rounds $<n$ (counts, pseudo-means, divergences, leader, reference, test divergences, threshold,
paced targets) are those computed from the history of the rounds $<r$; the count $N_c(r)$ is the
number of pulls of $c$ in the rounds $<r$. Along almost every run, the action of every round is
HY's rule applied to the history of the previous rounds; at a round of the IMED phase of batch
$\ell+1$ the pulled arm minimises the index $N_c(t-1)\tilde K_\ell(c)+y_\ell\log N_c(t-1)$ with
$\tilde K_\ell$ frozen at $t_\ell$, and this index is non-decreasing in $t$.
\end{lemma}

\begin{proof}
\uses{def:hy, def:learner}
Every statistic at the cutoff $r$ is a function of the actions and rewards of the rounds $<r$;
the actions of a deterministic learner are almost surely its policy applied to the history. The
index is non-decreasing since $\tilde K_\ell(c)\ge0$, $y_\ell\ge0$ and the counts grow.
\end{proof}

\begin{lemma}[the leader when $x>1$, on a run]\label{flem:leaderA-run}
\lean{Incognita.BatchedBandits.hyAlg_ae_integral_leader_eq_iSup, Incognita.BatchedBandits.hyAlg_ae_pseudoMean_reference_le, Incognita.BatchedBandits.HY.reference_spec, Incognita.BatchedBandits.HY.pseudoMean_le_leader, Incognita.BatchedBandits.HY.ofReal_le_mul_pseudoKinf_zero_add}
\uses{lem:structA, lem:thrA, lem:pseudo, lem:kinf, def:hy, def:model, def:constA, def:ranksA, def:evA}
Let $\ell\ge3$ be such that (R2) of Definition~\ref{def:ranksA} holds at $\ell-1$ and $\ell$,
with $\omega_{\ell-2},\omega_{\ell-1}\ge0$; let $\varepsilon>0$, an optimal arm $a^\star$ and
a real $\bar m$ with $\mathbf m_c\le\bar m$ for every sub-optimal $c$, such that
$\mathbf m^\star+\varepsilon<1$, $\bar m+\varepsilon<\mathbf m^\star-2\varepsilon$ and
$\mathbf m_c+\varepsilon<\mathbf m^\star-3\varepsilon$ for every sub-optimal $c$; let
$2w_0\le\varepsilon U_\ell$, and $z$ with
$z+w_0\kappa(\mathbf m^\star+\varepsilon)\le(1-\gamma_{\ell-1})y_{\ell-2}\mathsf L_{\ell-1}$.
Along almost every run of HY on $\nu\in\cP$, assume:
(A) for $j\in\{\ell-1,\ell\}$, every arm $c$ with $N_c(t_j)\ge\frac12U_j$ has
$|\hat{\mathbf m}_c(t_j)-\mathbf m_c|\le\varepsilon$;
(V) it is not the case that $N_{a^\star}(t_\ell)<\frac12U_\ell$ and
$N_{a^\star}(t_{\ell-1})\Kinf(\hat\nu_{a^\star}(t_{\ell-1}),\bar m+\varepsilon)\ge z$;
(P) no sub-optimal $c$ has $N_c(t_\ell)<\frac12U_\ell$,
$N_c(t_{\ell-1})\Kinf(\hat\nu_c(t_{\ell-1}),\mathbf m^\star+\varepsilon)\ge z$ and
$\hat{\mathbf m}_c(t_\ell)\ge\mathbf m^\star-3\varepsilon$
($\hat\nu_c(t_j)$, $\hat{\mathbf m}_c(t_j)$ the empirical law and mean of the rewards of $c$
before $t_j$). Then (i)--(iii) of Lemma~\ref{lem:leaderA} hold, and
$\hat\mu_{\ell-1}\le\mathbf m^\star+\varepsilon$ holds under (A) at $j=\ell-1$ alone (with
(R2) at $\ell-1$ only).
\end{lemma}

\begin{proof}
\uses{lem:structA, lem:thrA, lem:pseudo, lem:kinf, def:hy, def:model, def:evA}
The paper's proof of Lemma~\ref{lem:leaderA}, with the events of Definition~\ref{def:evA} read on
the run: on a run, $\hat\nu_c(t_j)$ is $\hat\nu_{c,N_c(t_j)}$ of the stack of $c$, with
$1\le N_c(t_{\ell-1})\le N_c(t_\ell)$, so an outcome of
$\mathcal A_{\ell-1}\cap\mathcal A_\ell\cap\mathcal V_\ell^c\cap\Pi_\ell^c$ satisfies (A), (V)
and (P) (with $\bar m=\mathbf m_2$ and the paper's $z_\ell$). The reference maximises the
pseudo-mean over the well-sampled arms and is well sampled when one is; the leader maximises it
over all arms; a rejected arm satisfies (eq:rejA) of the paper's proof for every level in
$[\hat\mu_{\ell-1},1)$.
\end{proof}

\begin{lemma}[a good batch when $x>1$, on a run]\label{flem:goodA-run}
\lean{Incognita.BatchedBandits.hyAlg_ae_le_mul_testKinf_and_pullCount_le_max, Incognita.BatchedBandits.HY.testKinf_le_imedKinf, Incognita.BatchedBandits.HY.target_le_of_rejected, Incognita.BatchedBandits.hyAlg_ae_pullCount_le_max'}
\uses{lem:structA, lem:thrA, lem:imedround, lem:pseudo, lem:kinf, def:hy, def:model, def:constA, def:ranksA}
Let $\ell\ge2$ be such that (R2) holds at $\ell$ and $\ell+1$, with $\omega_{\ell-1},\omega_\ell\ge0$,
$\gamma_\ell\ge0$; let $a$ be sub-optimal, $\varepsilon,D>0$, $K^+\ge0$, $x_+$ with
$K^++2\varepsilon\le(1+\omega_{\ell-1})D$, $D\mathsf L_{\ell-1}^{\beta-1}\ge2x_+^3$,
$x_{\ell-1},x_\ell\le x_+$, $0\le\mathbf m^\star-2\varepsilon$ and $\mathbf m^\star+\varepsilon<1$.
Along almost every run of HY on $\nu\in\cP$, assume the facts of Step~1 of the paper's proof of
Lemma~\ref{lem:goodA}: $\hat c_\ell$ is optimal,
$\mathbf m^\star-2\varepsilon\le\hat\mu_\ell\le\mathbf m^\star+\varepsilon$,
$\hat\mu_{\ell-1}\le\mathbf m^\star+\varepsilon$,
$\Kinf(\hat\nu_a(t_\ell),\mathbf m^\star-2\varepsilon)\ge D$,
$\Kinf(\hat\nu_a(t_{\ell-1}),\mathbf m^\star+\varepsilon)\le K^++\varepsilon$ and
$w_0\kappa(\mathbf m^\star+\varepsilon)\le\varepsilon N_a(t_{\ell-1})$. Then $a$ is rejected at
$t_\ell$, $\hat K_\ell(a)\ge D$, and every round $t\in(t_\ell,t_{\ell+1}]$ has
$N_a(t)\le\max(N_a(t_\ell),K+(1+\omega_\ell)y_\ell\log t/D)$.
\end{lemma}

\begin{proof}
\uses{flem:leaderA-run, lem:countsA, lem:structA, lem:thrA, lem:imedround, lem:pseudo, lem:kinf,
  def:hy, def:model, def:evA, def:constA, def:ranksA}
Steps~2--4 of the paper's proof of Lemma~\ref{lem:goodA}. With $D=\Delta_\varepsilon$ and
$K^+=\Kinf(\nu_a,\mathbf m^\star+\varepsilon)$, the hypotheses are Step~1 of the paper: they
follow, on $\mathcal G_\ell(\varepsilon)$, from Lemma~\ref{flem:leaderA-run} at $\ell-1$ and
$\ell$, Lemma~\ref{lem:countsA} and the traces of $\mathcal L_\ell^c$, $\mathcal U_\ell^c$ at the
counts $N_a(t_\ell),N_a(t_{\ell-1})\ge u_\ell(\varepsilon)$, with (T3). In the IMED phase,
Lemma~\ref{lem:imedround} is used with $\tilde K_\ell(a)\ge\hat K_\ell(a)\ge D$ in
$[0,\infty]$ ($\tilde K_\ell(a)=\infty$ is allowed: an arm pulled at an IMED round has a finite
divergence); in the paced phase the target of a rejected arm with $\hat K_\ell(a)\ge D$ is at
most $K+(1+\omega_\ell)y_\ell\log t/D$.
\end{proof}

\begin{lemma}[the pieces of the decomposition when $x>1$]\label{flem:decompA-run}
\lean{Incognita.BatchedBandits.hyAlg_ae_net_le, Incognita.BatchedBandits.pullCount_le_add_sub, Incognita.BatchedBandits.le_max_add_sum_of_step}
\uses{lem:structA, lem:imedround, def:hy, def:ranksA}
(Case~2 of (a).) Let (R2) hold at $\ell+1$, $\omega_\ell\ge0$, $x_\ell\le x_+$ and
$K+1+8x_+^2\mathsf L_{\ell+1}/\varepsilon'^2\le U_{\ell+1}$. Along almost every run of HY on
$\nu\in\cP$, an arm $a\ne\hat c_\ell$ with $\tilde K_\ell(a)\ge\varepsilon'^2/8$ has
$N_a(t)\le\max(N_a(t_\ell),2U_{\ell+1})$ for $t\in[t_\ell,t_{\ell+1}]$.
(Case~3 of (a).) For any sequence of actions, $N_c(t)\le N_c(s)+(t-s)$ for $s\le t$.
((b).) A real sequence with $N_{k+1}\le\max(N_k,\bar M,Y)+S_k$ and $S_k\ge0$ for $i_0\le k<j$ has
$N_k\le\max(N_{i_0},\bar M,Y)+\sum_{i_0\le i<k}S_i$ for $i_0\le k\le j$.
\end{lemma}

\begin{proof}
\uses{lem:structA, lem:imedround, def:hy}
Cases~2 and~3 and the induction of the paper's proof of Lemma~\ref{lem:decompA}, with a bound $Y$
of the nets in place of their running maximum.
\end{proof}

\begin{lemma}[a good batch when $x>1$, from the good event on a run]\label{flem:goodA-good}
\lean{Incognita.BatchedBandits.hyAlg_ae_pullCount_le_max_of_good}
\uses{flem:leaderA-run, flem:goodA-run, lem:countsA, def:hy, def:model, def:constA, def:ranksA,
  def:evA}
Let $\ell\ge4$ be such that (R2) holds at $\ell-2,\dots,\ell+1$, with $\omega\ge0$ there and
$\gamma_\ell\ge0$; let $a^\star$ be optimal, $a$ sub-optimal, $\bar m$ a bound of the sub-optimal
means and $\varepsilon>0$ with the conditions of Lemma~\ref{lem:constA}(a); let $2w_0\le\varepsilon
U_\ell$, $z+w_0\kappa(\mathbf m^\star+\varepsilon)\le(1-\gamma_{\ell-1})y_{\ell-2}\mathsf
L_{\ell-1}$, (T1) $\kappa_\varepsilon\mathsf L_j^{\beta-1}\ge1$ for $j=\ell-2,\ell-1$, a count
threshold $u\le\mathsf L_{\ell-2}/\kappa_\varepsilon$ with $w_0\kappa_\varepsilon\le\varepsilon u$,
and $D>0$, $K^+\ge0$, $x_+$ with $K^++2\varepsilon\le(1+\omega_{\ell-1})D$,
$D\mathsf L_{\ell-1}^{\beta-1}\ge2x_+^3$, $x_{\ell-1},x_\ell\le x_+$. Along almost every run of HY
on $\nu\in\cP$, assume the trace of $\mathcal G_\ell(\varepsilon)$ on the run: (A) at
$t_{\ell-2},t_{\ell-1},t_\ell$; (V) and (P) of Lemma~\ref{flem:leaderA-run} at the test $t_\ell$;
not ($N_a(t_\ell)\ge u$ and $\Kinf(\hat\nu_a(t_\ell),\mathbf m^\star-2\varepsilon)<D$); not
($N_a(t_{\ell-1})\ge u$ and $\Kinf(\hat\nu_a(t_{\ell-1}),\mathbf m^\star+\varepsilon)>K^++
\varepsilon$). Then every round $t\in(t_\ell,t_{\ell+1}]$ has
$N_a(t)\le\max(N_a(t_\ell),K+(1+\omega_\ell)y_\ell\log t/D)$.
\end{lemma}

\begin{proof}
\uses{flem:leaderA-run, flem:goodA-run, lem:countsA, lem:constA, def:evA}
Step~1 of the paper's proof of Lemma~\ref{lem:goodA}: Lemma~\ref{flem:leaderA-run} at the test
$t_\ell$ gives the optimal leader and the test levels; Lemma~\ref{lem:countsA} (along a run) at
$t_{\ell-1}$ and $t_\ell$ gives $N_a(t_{\ell-1}),N_a(t_\ell)\ge\mathsf
L_{\ell-2}/\kappa_\varepsilon\ge u$, so the traces of $\mathcal L_\ell^c$, $\mathcal U_\ell^c$
apply, and (T3) holds at the count $N_a(t_{\ell-1})$; $\mathbf m^\star-2\varepsilon>\bar m+
\varepsilon\ge\mathbf m_a\ge0$. Steps~2--4 are Lemma~\ref{flem:goodA-run}. The events
$\mathcal V_{\ell-1}^c$, $\Pi_{\ell-1}^c$ of $\mathcal G_\ell$ serve only the claim that
$\hat c_{\ell-1}$ is optimal, which the bound does not use: they are not assumed.
\end{proof}

\begin{lemma}[the inequality (eq:DA) on a run]\label{flem:decompA-DA}
\lean{Incognita.BatchedBandits.hyAlg_ae_pullCount_le_max_add}
\uses{flem:goodA-good, flem:leaderA-run, flem:decompA-run, def:hy, def:model, def:constA,
  def:ranksA, def:evA}
Under the assumptions of Lemma~\ref{flem:goodA-good} at the accuracy $\varepsilon$, let
$\varepsilon'>0$ satisfy the conditions of Lemma~\ref{lem:constA}(a), $2w_0\le\varepsilon'U_\ell$,
$z+w_0\kappa(\mathbf m^\star+\varepsilon')\le(1-\gamma_{\ell-1})y_{\ell-2}\mathsf L_{\ell-1}$ and
(T4) $K+1+8x_+^2\mathsf L_{\ell+1}/\varepsilon'^2\le U_{\ell+1}$. Let $G$ be the trace of
$\mathcal G_\ell(\varepsilon)$ on the run (the assumptions on the run of
Lemma~\ref{flem:goodA-good}) and $F^c$ that of $\mathcal F_\ell(\varepsilon')^c$ ((A) at
$t_{\ell-1},t_\ell$, (V) and (P) at the test $t_\ell$, at the accuracy $\varepsilon'$). Along almost
every run of HY on $\nu\in\cP$, every round $t\in(t_\ell,t_{\ell+1}]$ has
\[
N_a(t)\le\max\Big(N_a(t_\ell),\ K+\frac{(1+\omega_\ell)y_\ell\log t}{D},\ 2U_{\ell+1}\mathbf
1(G^c)\Big)+(t-t_\ell)\mathbf 1(F).
\]
\end{lemma}

\begin{proof}
\uses{flem:goodA-good, flem:leaderA-run, flem:decompA-run}
The three cases of the paper's proof of Lemma~\ref{lem:decompA}(a): on $G$,
Lemma~\ref{flem:goodA-good}; on $G^c\cap F^c$, Lemma~\ref{flem:leaderA-run} at the accuracy
$\varepsilon'$ gives an optimal leader and $\tilde K_\ell(a)\ge\hat K_\ell(a)\ge\varepsilon'^2/8$,
and Case~2 of Lemma~\ref{flem:decompA-run} gives $N_a(t)\le\max(N_a(t_\ell),2U_{\ell+1})$; on $F$,
Case~3.
\end{proof}

\begin{lemma}[the failure event when $x>1$, on a run]\label{flem:failA}
\lean{Incognita.BatchedBandits.measure_not_accurate_and_not_rejected_le}
\uses{lem:probasA, lem:constA, def:evA, def:model, def:hy}
For every learner on $\nu\in\cP$, a test at $t_\ell$, $\varepsilon>0$ with $\mathbf
m^\star+\varepsilon<1$, an optimal $a^\star$, $\bar m$ with $\bar m+\varepsilon\ge0$ and
$2\eta(\bar m+\varepsilon)\le\mathbf m^\star-(\bar m+\varepsilon)$ ($\eta\ge0$), tilts
$\theta_c\ge0$ of the sub-optimal arms with the moment conditions of Lemma~\ref{lem:pair} at the
level $\mathbf m^\star+\varepsilon$ and the mean level $\mathbf m^\star-3\varepsilon$, and $z>0$: the
probability of the trace $F$ of $\mathcal F_\ell(\varepsilon)$ on the run
(Lemma~\ref{flem:decompA-DA}) is at most
$Kp_{\mathcal A}(U_{\ell-1})+Kp_{\mathcal A}(U_\ell)+C(\bar m+\varepsilon)U_\ell^3e^{-(1+2\eta)z}+
KC(\mathbf m^\star+\varepsilon)U_\ell^4e^{-(1+2\eta)z}$, with
$p_{\mathcal A}(U)=2e^{-\varepsilon^2U}/(1-e^{-2\varepsilon^2})$ and
$C(r)=(4+2\kappa(r))e^{2(1+2\eta)}$.
\end{lemma}

\begin{proof}
\uses{lem:probasA, def:evA}
$F$ is contained in the union of the traces of $\mathcal A_{\ell-1}^c$, $\mathcal A_\ell^c$,
$\mathcal V_\ell$ and of the $\Pi_\ell$ of each sub-optimal arm; Lemma~\ref{lem:probasA}(a)--(b)
(for the traces, every learner) bounds each, and a union bound over at most $K$ sub-optimal arms.
\end{proof}

\begin{lemma}[the lost batches are summable, on a run]\label{flem:failA-sum}
\lean{Incognita.BatchedBandits.sum_batch_mul_measure_not_accurate_and_not_rejected_le}
\uses{flem:failA, lem:sumsA, lem:ranksA, def:evA, def:ranksA, def:grid}
Under the assumptions of Lemma~\ref{flem:failA} with $\eta>0$ and $\beta\ge1$, let levels $z_\ell>0$
satisfy (R3) $(1+2\eta)z_\ell\ge(1+\eta)\mathsf L_{\ell+1}$ and $\mathsf L_{\ell+1}\le x_+^3\mathsf
L_{\ell-2}$ for $\ell\in[k_0,L)$ ($k_0\ge3$), and $S\ge2x_+^3y-\varepsilon^2y^\beta$ for every
$y\ge0$. Then
$\sum_{k_0\le\ell<L}B_{\ell+1}\PP(F_\ell)\le\frac c2e^S+(1+K)C(\mathbf
m^\star+\varepsilon)\sigma_{4\beta}(\eta)$ with $c=4K/(1-e^{-2\varepsilon^2})$, where $F_\ell$ is
the trace of Lemma~\ref{flem:failA} at the level $z_\ell$.
\end{lemma}

\begin{proof}
\uses{flem:failA, lem:sumsA}
Lemma~\ref{flem:failA} gives $\PP(F_\ell)\le a_\ell+v_\ell$ with
$a_\ell=ce^{-\varepsilon^2\mathsf L_{\ell-2}^\beta}$ ($U_{\ell-1}\ge\mathsf L_{\ell-2}^\beta$,
$U_\ell\ge\mathsf L_{\ell-1}^\beta\ge\mathsf L_{\ell-2}^\beta$) and $v_\ell=(1+K)C(\mathbf
m^\star+\varepsilon)U_\ell^4t_{\ell+1}^{-(1+\eta)}$ ($\kappa(\bar m+\varepsilon)\le\kappa(\mathbf
m^\star+\varepsilon)$, $U^3\le U^4$, and $e^{-(1+2\eta)z_\ell}\le t_{\ell+1}^{-(1+\eta)}$ by (R3));
Lemma~\ref{lem:sumsA}(b) in its general form sums them.
\end{proof}

\begin{lemma}[good high batches when $x=1$, on a run]\label{flem:goodhigh-run}
\lean{Incognita.BatchedBandits.hyAlg_ae_le_log_pullCount_leader, Incognita.BatchedBandits.hyAlg_ae_card_le_card_filter_le_pullCount_leader_add, Incognita.BatchedBandits.card_filter_lt_le_card}
\uses{def:hy, def:high, def:water, def:model, def:learner, def:grid}
Call a batch $i\ge2$ \emph{high at the level $v$ for the arm $c$} if
$J_{i-1}(c,t_i+1)\ge y_{i-1}v$ (by (W1), a $(\delta,j)$-high batch is high at the level
$(1-\delta)\mathsf L_j$ for every arm). Along almost every run of HY ($w_0\ge0$) on $\nu\in\cP$, let
$S$ be a finite set of batches $i\ge2$, each high at the level $v$ for its leader $\hat c_{i-1}$,
and $u\le e^v$. Then $N_{\hat c_{i-1}}(t_i)\ge u$ for every $i\in S$, and at most $K$ batches
$i\in S$ have $N_{\hat c_{i-1}}(t_{i-1})<u$.
\end{lemma}

\begin{proof}
\uses{lem:imedround, def:hy, def:high, def:water}
Steps~1--2 of the paper's proof of Lemma~\ref{lem:goodhigh}. The leader has $\tilde K_{i-1}(\hat c_{i-1})=0$
(Step~1 of the proof of Lemma~\ref{lem:imedround}), so its index is $y_{i-1}\log N_{\hat c_{i-1}}(t_i)$,
and $y_{i-1}\ge1$; then the injectivity of Step~2.
\end{proof}

\begin{lemma}[few observations when $x=1$, on a run]\label{flem:fewC-run}
\lean{Incognita.BatchedBandits.hyAlg_ae_exists_mem_pullCount_lt_and_le_mul_pseudoKinf, Incognita.BatchedBandits.hyAlg_ae_exists_mem_pullCount_lt_and_le_mul_pseudoKinf_of_pseudoMean_leader_lt, Incognita.BatchedBandits.hyAlg_ae_pullCount_lt_and_exists_mem_le_mul_pseudoKinf_of_le_pseudoMean, Incognita.BatchedBandits.ofReal_sub_log_le_div, Incognita.BatchedBandits.HY.sub_two_mul_le_pseudoMean, Incognita.BatchedBandits.HY.pseudoMean_le_add_of_abs_sub_le, Incognita.BatchedBandits.HY.sub_two_mul_le_pseudoMean_leader}
\uses{def:hy, def:high, def:evC, def:model, def:kinf, def:stacks, def:grid}
Along almost every run of HY ($w_0\ge0$) on $\nu\in\cP$:
(common step) let $c$ be an arm, $S$ a nonempty set of batches $i\in[2,J]$, high at the level $v$
for $c$, with $\tilde{\mathbf m}^\star_{i-1}\le r'<1$, and let $N_c(t_J)<u$, $Y=v-\log u>0$, $\delta<1$,
$u\le(1-\delta)^{-(|S|-1)}Y/\kappa(r')$. Then some $i\in S$ has $n=N_c(t_{i-1})\in[1,u)$ and
$n\Kinf(\hat\nu_c(t_{i-1}),r')\ge(1-\delta)Y-w_0\kappa(r')$.
(First inclusion) Let $a^\star$ be optimal, $\mathbf m_c\le\mathbf m_2$ for the sub-optimal $c$, $\varepsilon\ge0$,
$\mathbf m_2+\varepsilon<\mathbf m^\star-2\varepsilon$, $\mathbf m^\star+\varepsilon<1$, $w_0\le\varepsilon u$; let $S$ be as above for
$c=a^\star$, with $N_{\hat c_{i-1}}(t_{i-1})\ge u$ for $i\in S$, and the size condition at $r'=\mathbf m_2+\varepsilon$.
If every arm with $N_c(t_k)\ge u$, $k\le J$, has $|\hat{\mathbf m}_c(t_k)-\mathbf m_c|\le\varepsilon$ (the trace of
$\mathcal D_j^c$) and $\tilde{\mathbf m}^\star_J<\mathbf m^\star-2\varepsilon$ (the class $\mathcal U_j$), the conclusion of the common
step holds for $c=a^\star$, $r'=\mathbf m_2+\varepsilon$ (the trace of $\mathcal V^\star_j$).
(Second inclusion) With $\mathbf m_a+\varepsilon<\mathbf m^\star-3\varepsilon$, the same $S$ for $c=a$ and the size condition at
$r'=\mathbf m^\star+\varepsilon$, the trace of $\mathcal D_j^c$ and $\hat{\mathbf m}_a(t_J)\ge\mathbf m^\star-3\varepsilon$ give $N_a(t_J)<u$ and the
conclusion of the common step for $c=a$, $r'=\mathbf m^\star+\varepsilon$ (the trace of $\mathcal V_{a,j}$).
\end{lemma}

\begin{proof}
\uses{flem:goodhigh-run, lem:highC, lem:chain, lem:pseudo, lem:kinf, def:hy, def:evC, def:model, def:stacks}
The paper's proof of Lemma~\ref{lem:fewC}, with $J=j-1$, $S=S_{\rm go}$; on a run,
$\hat\nu_c(t_k)$ is $\hat\nu_{c,N_c(t_k)}$ of the stack of $c$, so the paper's events contain their
traces. The arithmetic $(1-\delta)(v-\log u)-w_0\kappa\ge(1-4\delta)\mathsf L_j$ for
$v=(1-\delta)\mathsf L_j$, $u=u_j$ is Lemma~\ref{lem:constC}(a), and the size condition is
Lemma~\ref{lem:goodhigh} with Lemma~\ref{lem:constC}(a).
\end{proof}

\begin{lemma}[an optimal arm in $W$ when $x=1$, on a run]\label{flem:optW-run}
\lean{Incognita.BatchedBandits.hyAlg_ae_integral_reference_eq_iSup, Incognita.BatchedBandits.hyAlg_ae_exists_integral_eq_iSup_and_le_two_mul_pullCount}
\uses{def:hy, def:high, def:evC, def:model, def:kinf, def:stacks, def:grid}
Along almost every run of HY ($w_0\ge0$) on $\nu\in\cP$, with $a^\star$, $\mathbf m_2$, $\varepsilon$ as in the first
inclusion of Lemma~\ref{flem:fewC-run}, $U\ge1$, $2w_0\le\varepsilon U$, $2u\le U\le2e^v$, a set $S$ of at
least two batches $i\in[2,J]$, high at the level $v$ for every arm, with
$N_{\hat c_{i-1}}(t_{i-1})\ge u$, and $u\le(1-\delta)^{-(|S|-2)}(v-\log u)/\kappa(\mathbf m_2+\varepsilon)$: if the traces
of $\mathcal D_j^c$ and $(\mathcal V^\star_j)^c$ up to $t_J$ hold, then some optimal arm has $N(t_J)\ge U/2$, the
reference $r_J$ (net size $U$) is optimal, $\mathbf m^\star-2\varepsilon\le\hat\mu_J\le\mathbf m^\star+\varepsilon$, and every sub-optimal
$c$ with $N_c(t_J)\ge U/2$ has $\tilde{\mathbf m}_c(t_J)<\mathbf m^\star-2\varepsilon$.
\end{lemma}

\begin{proof}
\uses{flem:goodhigh-run, flem:fewC-run, lem:pseudo, def:hy, def:evC, def:model}
The paper's proof of Lemma~\ref{lem:optW}: Step~1 by Lemma~\ref{flem:goodhigh-run}, Step~3 by the
common step of Lemma~\ref{flem:fewC-run} on $S$ without its largest batch, Step~4 by the
definition of the reference.
\end{proof}

\begin{lemma}[pulls of $a$ when $x=1$, on a run]\label{flem:pullsC-run}
\lean{Incognita.BatchedBandits.HY.pseudoMean_lt_and_ne_leader_and_pseudoKinf_le_testKinf, Incognita.BatchedBandits.hyAlg_ae_pullCount_le_card_add_div, Incognita.BatchedBandits.hyAlg_ae_pullCount_le_max_of_le_imedKinf, Incognita.BatchedBandits.pullCount_le_max_of_forall_add_one, Incognita.BatchedBandits.hyAlg_ae_pullCount_add_one_le}
\uses{def:hy, def:hy-vocab, def:model, def:kinf, def:grid}
(i) For one history with rewards in $[0,1]$, $N_a\ge1$, $0\le\mu\le\hat\mu<1$ and
$D=\Kinf(\hat\nu_a,\mu)>0$: $\hat{\mathbf m}_a<\mu$, $a$ is not the leader and $\tilde K(a)\ge\hat K(a)\ge D$.
Along almost every run of HY ($w_0\ge0$) on $\nu\in\cP$, for $\ell\ge1$ and $a\ne\hat c_\ell$:
(ii) if $\omega_\ell\ge0$, $\hat K_\ell(a)\ge D>0$, the batch $\ell+1$ is not activated or $a$ is
rejected at $t_\ell$, and $a$ is pulled in the batch, then
$N_a(t_{\ell+1})\le K+(1+\omega_\ell)y_\ell\mathsf L_{\ell+1}/D$;
(iii) if $\vartheta_\ell\ge0$ and $\tilde K_\ell(a)\ge D>0$, then
$N_a(t_{\ell+1})\le\max(N_a(t_\ell),1+y_\ell\log s_\ell/D,K-1+U_{\ell+1})$.
\end{lemma}

\begin{proof}
\uses{lem:imedround, lem:service, lem:kinf, lem:pseudo, def:hy, def:hy-vocab}
The paper's proof of Lemma~\ref{lem:pullsC}. Every pull of $a$ at a round $k$ of the batch has
$N_a(k)+1\le g(k)$ for a nondecreasing $g$: at an IMED round $g(k)=1+y_\ell\log(k+1)/D$
(Lemma~\ref{lem:imedround}), at a paced round the target (the paced rule), at most
$K+(1+\omega_\ell)y_\ell\log(k+1)/D$ for a rejected arm and $K-1+U_{\ell+1}$ in general; then
$N_a(t_{\ell+1})\le\max(N_a(t_\ell),g(t_{\ell+1}-1))$, and $N_a(t_{\ell+1})\le g(t_{\ell+1}-1)$ if $a$ is
pulled (the last pull). With $D\ge\varepsilon^2/8$, $\log s_\ell\le2\mathsf L_\ell$ and $y_\ell\le\bar y$, (iii)
gives the paper's $N_a(t_{\ell+1})-N_a(t_\ell)\le c_\ell$.
\end{proof}

\begin{lemma}[logarithmic count when $x=1$, on a run]\label{flem:logcountC-run}
\lean{Incognita.BatchedBandits.hyAlg_ae_le_mul_pullCount_of_le_imedIndex_leader, Incognita.BatchedBandits.hyAlg_ae_le_mul_pullCount_of_le_imedIndex}
\uses{def:hy, def:high, def:evC, def:model, def:kinf, def:grid}
Along almost every run of HY ($w_0\ge0$) on $\nu\in\cP$, let $\varepsilon\ge0$ with $\mathbf m^\star+\varepsilon<1$ and let
$i\in[2,J]$ be a batch high at the level $v$ for an arm $c$, whose leader has
$N_{\hat c_{i-1}}(t_{i-1})\ge u$, such that every arm with $N_{c'}(t_{i-1})\ge u$ has
$\hat{\mathbf m}_{c'}(t_{i-1})\le\mathbf m_{c'}+\varepsilon$. Then $v\le(\kappa(\mathbf m^\star+\varepsilon)+1)N_c(t_J)$.
\end{lemma}

\begin{proof}
\uses{lem:kinf, lem:pseudo, def:hy, def:high}
The paper's proof of Lemma~\ref{lem:logcountC}: $\tilde{\mathbf m}^\star_{i-1}\le\hat{\mathbf m}_{\hat c_{i-1}}(t_{i-1})\le
\mathbf m^\star+\varepsilon$, so $\tilde K_{i-1}(c)\le\kappa(\mathbf m^\star+\varepsilon)$; the index at the end of the batch is at
most $\kappa N_c(t_i)+y_{i-1}\log N_c(t_i)$, $\log N\le N$, and $N_c(t_i)\le N_c(t_J)$.
\end{proof}

\begin{lemma}[the bad events when $x=1$, on a run]\label{flem:controlC-run}
\lean{Incognita.BatchedBandits.measure_exists_mem_le_pullCount_abs_sub_gt_le, Incognita.BatchedBandits.measure_exists_mem_lt_mul_pseudoKinf_le, Incognita.BatchedBandits.measure_exists_mem_lt_mul_pseudoKinf_and_le_le, Incognita.BatchedBandits.lintegral_sum_ite_pseudoKinf_lt_le, Incognita.BatchedBandits.sum_Icc_mul_exp_neg_le, Incognita.BatchedBandits.sum_log_rewardByCount_le_of_pseudoKinf_lt}
\uses{def:evC, def:stacks, def:model, def:kinf, def:learner}
For every learner on $\nu\in\cP$ and every set $T$ of cutoffs:
(a) for $u,\varepsilon>0$, the probability that some arm $c$ and some $s\in T$ have $N_c(s)\ge u$ and
$|\hat{\mathbf m}_c(s)-\mathbf m_c|>\varepsilon$ is at most $K\cdot2e^{-2\varepsilon^2u}/(1-e^{-2\varepsilon^2})$;
(b) for an arm $a'$, $r\in[0,1)$, $\eta\ge0$ with $2\eta r\le\mathbf m_{a'}-r$, $z>0$ and $U$, the probability that
some $s\in T$ has $N_{a'}(s)<U$ and $N_{a'}(s)\Kinf(\hat\nu_{a'}(s),r)\ge z$ is at most
$C_\eta U^3e^{-(1+2\eta)z}$, $C_\eta=(4+2\kappa(r))e^{2(1+2\eta)}$;
(c) for an arm $c$, a cutoff $t$, $r\in[0,1)$, $q$, $\theta,\eta\ge0$ with the moment conditions of
Lemma~\ref{lem:pair} for $\nu_c$, $z>0$ and $U$, the probability that $N_c(t)<U$,
$\hat{\mathbf m}_c(t)\ge q$ and some $s\in T$, $s\le t$, has $N_c(s)\Kinf(\hat\nu_c(s),r)\ge z$ is at most
$C_\eta U^4e^{-(1+2\eta)z}$;
(d) for a grid, an arm $a$ with $N_a(t_1)\ge1$ almost surely, $\mu\in[0,1)$, $\varepsilon'>0$,
$\lambda\in[0,\frac1{1-\mu})$ and $A=\E_{\nu_a}[\log(1-\lambda(Y-\mu))]$,
$\E\big[\sum_{2\le j\le L}N_a(t_{j-1})\mathbf 1(\Kinf(\hat\nu_a(t_{j-1}),\mu)<A-3\varepsilon'/4,\ N_a(t_{j-1})<N_a(t_j))\big]
\le1/\rho'^2$ for every $L$, with $\rho'$ of Lemma~\ref{lem:lowerdev}.
\end{lemma}

\begin{proof}
\uses{lem:hoeffding, lem:single, lem:pair, lem:lowerdev, lem:distinct, lem:stacks, def:evC, def:stacks}
On the product with the reward stacks, $\hat\nu_c(s)$ is $\hat\nu_{c,N_c(s)}$ (Lemma~\ref{lem:stacks}), so each
event is contained in the paper's event $\mathcal D_j$, $\mathcal V^\star_j$ or $\mathcal V_{a,j}$ on the stacks; then the
paper's proof of Lemma~\ref{lem:controlC}(a)--(c). For (d), the batches in which the count of $a$
increases start at distinct counts $n\in[1,t_L]$ (Lemma~\ref{lem:distinct}(i)); at the count $n$ the
event is in the event of Lemma~\ref{lem:lowerdev}, of probability at most $e^{-\rho'n}$, and
$\sum_{n\ge1}ne^{-\rho'n}\le1/\rho'^2$.
\end{proof}

\begin{lemma}[the consistency test when $x=1$, for one history]\label{flem:amb-hist}
\lean{Incognita.BatchedBandits.HY.ofReal_mul_threshold_le_count_mul_testKinf, Incognita.BatchedBandits.HY.testKinf_le_old, Incognita.BatchedBandits.HY.le_testKinf_new, Incognita.BatchedBandits.HY.ofReal_threshold_le_of_mem, Incognita.BatchedBandits.HY.threshold_le, Incognita.BatchedBandits.HY.threshold_nonneg}
\uses{def:hy, def:amb, def:kinf, def:model}
For a history of HY with rewards in $[0,1]$, cutoffs $r=t_\ell$, $r'=t_{\ell-1}$, net sizes $U=U_\ell$,
$U'=U_{\ell-1}$ and an arm $a$ with $N=N_a(t_\ell)$, let $y'=y_{\ell-1}\in(0,\bar y]$, $\gamma\in(0,\frac12]$,
$\mathsf L_\ell\ge0$, $\varepsilon>0$, $\Delta\ge0$. If $\hat K_\ell(a)\ge\varepsilon^2/8$, $a\ne\hat c_{\ell-1}$,
$\bar y\mathsf L_\ell\le(U/2)(\varepsilon^2/8)$, $2\bar y\log N\le\gamma(\varepsilon^2/8)N$,
$\hat K_{\ell-1}(a)\le\Delta(1+3\gamma/32)$ and $\hat K_\ell(a)\ge\Delta(1-\gamma/16)$, then $a$ is rejected at $t_\ell$:
$\hat K_\ell(a)>0$ and $N\hat K_\ell(a)\ge(1-\gamma)y'\hat\Lambda_\ell$. Moreover, with $m=\mathbf m^\star$, $\zeta>0$:
$\hat K_\ell(a)\ge\Delta-2\zeta'$ when $m-2\zeta\le\hat\mu_\ell<1$, $2\zeta\le m<1$,
$\Kinf(\hat\nu_a(t_\ell),m)\ge\Delta-\zeta'$ and $2\zeta/(1-m)\le\zeta'$; and $\hat K_{\ell-1}(a)\le\Delta+3\zeta'$ when
$\hat\mu_{\ell-1}\le m+\zeta<1$, $w_0\kappa(m+\zeta)\le\zeta'N_a(t_{\ell-1})$, $\Kinf(\hat\nu_a(t_{\ell-1}),m)\le\Delta+\zeta'$
and $\zeta/(1-m-\zeta)\le\zeta'$.
\end{lemma}

\begin{proof}
\uses{lem:kinf, lem:pseudo, def:hy, def:amb}
Steps~3--4 of the paper's proof of Lemma~\ref{lem:amb}: a well-sampled $a$ has
$N\hat K_\ell(a)\ge(U/2)\varepsilon^2/8\ge\bar y\mathsf L_\ell\ge y'\hat\Lambda_\ell$; otherwise $a\in Q_\ell$, so
$\hat\Lambda_\ell\le N\hat K_{\ell-1}(a)/y'+\log N$ (definition of $\hat\Lambda_\ell$ as a minimum), and the two cases
$\log N>\frac\gamma2\hat\Lambda_\ell$, $\log N\le\frac\gamma2\hat\Lambda_\ell$ conclude, with
$(1-\frac\gamma2)(1+\frac{3\gamma}{32})\le1-\frac\gamma{16}$. With $\zeta'=\gamma\Delta/32$ the two bounds of the
second part are the hypotheses of the first (Lemma~\ref{lem:kinf}(iii), (v) and
Lemma~\ref{lem:pseudo}(i)--(ii)).
\end{proof}

\begin{lemma}[the consistency test when $x=1$, on a run]\label{flem:amb-run}
\lean{Incognita.BatchedBandits.hyAlg_ae_ofReal_mul_threshold_le_count_mul_testKinf, Incognita.BatchedBandits.exists_le_pullCount_of_card_mul_le, Incognita.BatchedBandits.mul_log_le_mul_of_le}
\uses{def:hy, def:amb, def:kinf, def:model, def:grid}
Along almost every run of HY ($w_0\ge0$) on $\nu\in\cP$, let $\ell\ge2$, $a$ sub-optimal, $\varepsilon,\zeta>0$,
$\gamma=\gamma_\ell\in(0,\frac12]$, $\Delta\ge0$, $\zeta'=2\zeta/(1-\mathbf m^\star)=\gamma\Delta/32$, $\bar y\ge y_{\ell-1}$ and
$n_\ell\ge e$, with the conditions (b)--(d) of Definition~\ref{def:amb} (at the count $n_\ell$),
$t_{\ell-1}\ge KU_{\ell-1}$ and $U_{\ell-1}\le U_\ell$. Assume: some optimal arm has $N(t_\ell)\ge U_\ell/2$;
$\hat\mu_\ell\ge\mathbf m^\star-2\varepsilon$; $\hat{\mathbf m}_a(t_\ell)<\mathbf m^\star-3\varepsilon$; $N_a(t_\ell),N_a(t_{\ell-1})\ge n_\ell$;
$a\ne\hat c_{\ell-1}$; the traces of $\mathsf{Acc}_\ell$ (at $t_{\ell-1}$, $t_\ell$, for the arms with $N\ge U_{\ell-1}/2$) and
of $\mathsf{KC}_\ell$ (at $N_a(t_{\ell-1})$, $N_a(t_\ell)$). Then $a\ne\hat c_\ell$, $\hat K_\ell(a)\ge\varepsilon^2/8$, and $a$ is
rejected at $t_\ell$.
\end{lemma}

\begin{proof}
\uses{flem:amb-hist, flem:pullsC-run, lem:kinf, lem:pseudo, def:hy, def:amb}
Steps~1--4 of the paper's proof of Lemma~\ref{lem:amb}. The test levels: $\hat\mu_\ell\le\mathbf m^\star+\zeta$ and
$\hat\mu_\ell\ge\tilde{\mathbf m}_{a'}(t_\ell)\ge\mathbf m^\star-2\zeta$ for the well-sampled optimal $a'$; some arm is well
sampled at $t_{\ell-1}$ since $\sum_cN_c(t_{\ell-1})=t_{\ell-1}\ge KU_{\ell-1}$, so $\hat\mu_{\ell-1}\le\mathbf m^\star+\zeta$. Then
$\Kinf(\hat\nu_a(t_\ell),\mathbf m^\star-2\varepsilon)\ge\varepsilon^2/8$ (Lemma~\ref{lem:kinf}(iv)), Lemma~\ref{flem:pullsC-run}(i), the
condition (c) at $N_a(t_\ell)\ge n_\ell$ ($u\mapsto\log u/u$ is non-increasing on $[e,\infty)$), and
Lemma~\ref{flem:amb-hist}.
\end{proof}

\begin{lemma}[many good high batches when $x=1$, on a run]\label{flem:levelC-run}
\lean{Incognita.BatchedBandits.hyAlg_ae_le_card_filter_le_imedIndex, Incognita.BatchedBandits.hyAlg_ae_le_card_filter_le_imedIndex_and_le_pullCount_leader, Incognita.BatchedBandits.exists_action_eq_of_pullCount_lt}
\uses{def:hy, def:water, def:high, def:grid, def:ass-x, def:model}
On a grid with $x_\ell\to1$, let $g\in(0,1)$, $m'$, and a rank $\ell_W\ge2$ from which every batch is
not activated or has $s_\ell-t_\ell\ge Kt_\ell$. There is a rank $L_0$, independent of the run, such
that along almost every run of HY every $L\ge L_0$ has at least $m'$ batches $j\in[2,L)$ high at the
level $(1-g)\mathsf L_L$ for every arm. For HY with $w_0\ge0$ on $\nu\in\cP$ and every $M$, there is such
a rank after which, for every $u\le e^{(1-g)\mathsf L_L}$, at least $M+2$ of these batches have a leader
with $N_{\hat c_{j-1}}(t_{j-1})\ge u$.
\end{lemma}

\begin{proof}
\uses{lem:W, lem:lowmass, lem:levelC, flem:goodhigh-run, def:water, def:high}
Take the water level $\Lambda_j=\min(\min_cJ_{j-1}(c,t_j+1),\,y_{j-1}\log t_{j-1})/y_{j-1}$, finite and
non-negative. An arm pulled in batch $j$ is pulled at some round of it, so (W2) at that round
(Lemma~\ref{lem:W}(i)) and $N_c(t_{j-1})\le t_{j-1}$ give $\log N_c(t_{j-1})\le\Lambda_j$; then
Lemma~\ref{lem:lowmass} and the abstract Lemma~\ref{lem:levelC}, whose rank does not depend on
$\Lambda$. A batch with $\Lambda_j\ge(1-g)\mathsf L_L$ is high for every arm. The second claim is
Lemma~\ref{flem:goodhigh-run} with $m'=M+K+2$.
\end{proof}

\begin{lemma}[the decomposition when $x=1$, abstract form]\label{flem:decompC-abstract}
\lean{Incognita.BatchedBandits.le_add_sum_of_classes, Incognita.BatchedBandits.sum_Ioc_sub_pred, Incognita.BatchedBandits.le_add_div_add_mul_of_le_add_div, Incognita.BatchedBandits.le_mul_of_le_add_div_sq}
\uses{}
Let $(N_j)$ be a nondecreasing sequence of integers, $1\le j_0\le L$, and let every $j\in(j_0,L]$
be bad, good, other or ambiguous, with $n_j=N_j-N_{j-1}\le B_j$ if bad, $N_j\le M$ if good and
$n_j>0$, $N_j\le CN_{j-1}$ and $Q(N_{j-1})$ if other and $n_j>0$, and $n_j\le c_j$ if ambiguous
($B,c,M,C\ge0$). Then
\[
N_L\le N_{j_0}+M+\sum_{j_0<j\le L}B_j\mathbf 1(\text{bad}_j)
+C\sum_{2\le j\le L}N_{j-1}\mathbf 1(Q(N_{j-1}),\,n_j>0)+\sum_{j_0<j\le L}c_j\mathbf 1(\text{amb}_j).
\]
The bounds of the two classes: $N\le K+(1+\omega)y\mathsf L_j/D$ with $D\ge D'$, $(1+\omega)y\le1+\varepsilon$, $\mathsf L_j\le\mathsf L_L$ and
$(1+\varepsilon)/D'\le1/\Delta+C_1\varepsilon$ gives $N\le K+\mathsf L_L/\Delta+C_1\varepsilon\mathsf L_L$; $N\le K+16\bar y\mathsf L_j/\varepsilon^2$ with
$(1-\delta)\mathsf L_j\le(\kappa+1)N'$, $N'\ge1$, gives $N\le C_\varepsilon N'$.
\end{lemma}

\begin{proof}
\uses{lem:distinct}
Steps~1--4 of the paper's proof of Lemma~\ref{lem:decompC}: telescoping, one bound per class, and
Lemma~\ref{lem:distinct}(ii) for the good increments.
\end{proof}

\begin{lemma}[shifting the level of $\Kinf$, primal form]\label{flem:kinf-shift}
\lean{Incognita.BatchedBandits.dInf_unitLaws_le_add, Incognita.BatchedBandits.dInf_unitLaws_sub_le_of_le, Incognita.BatchedBandits.dInf_mono, Incognita.BatchedBandits.dInf_unitLaws_ne_top, Incognita.klDiv_smul_add_smul_le_add, Incognita.klDiv_le_add_ofReal_neg_log_of_smul_le}
\uses{def:kinf}
For probability measures $P,Q,Q'$ and $c>0$ with $cQ\le Q'$, $\KL(P,Q')\le\KL(P,Q)-\log c$; in particular, for
$\rho$ a probability measure and $u\in[0,1)$, $\KL(P,(1-u)Q+u\rho)\le\KL(P,Q)-\log(1-u)$. Hence, for a law $P$ on
$[0,1]$, $\varepsilon>0$ and $\mu<1$ (the paper's $\varepsilon\le\mu$ is not needed: $\Kinf$ is defined, and $0$, below the
mean of $P$),
$\Kinf(P,\mu)\le\Kinf(P,\mu-\varepsilon)+\varepsilon/(1-\mu)$ (Lemma~\ref{lem:kinf}(v) for every law, with the primal
formula of the frozen statements), and $\Kinf(P,\cdot)$ is non-decreasing. Moreover $\Kinf(P,\mu)<\infty$ for every $\mu<1$.
\end{lemma}

\begin{proof}
\uses{def:kinf, lem:kinf}
Finiteness: if $\mu$ is below the mean $p$ of $P$, $P$ is a competitor and $\Kinf(P,\mu)=0$; otherwise the competitors
$(1-u)P+u\delta_1$ bound it by $\log\frac{1-p}{1-\mu}\le\kappa(\mu)$, as in Lemma~\ref{lem:kinf}(i). The mixture dominates $(1-u)Q$, so $dQ/d((1-u)Q+u\rho)\le1/(1-u)$ and, by the chain rule,
$\log\frac{dP}{d((1-u)Q+u\rho)}\le\log\frac{dP}{dQ}-\log(1-u)$ $P$-almost everywhere. For the shift, a competitor
$Q$ at the level $\mu-\varepsilon$ gives the competitor $(1-v)Q+v\delta_1$ at the level $\mu$,
$v=\varepsilon/(1-\mu+\varepsilon)$, and $-\log(1-v)=\log(1+\frac{\varepsilon}{1-\mu})\le\frac{\varepsilon}{1-\mu}$.
\end{proof}

\begin{lemma}[the final constant when $x=1$]\label{flem:constC-b}
\lean{Incognita.BatchedBandits.dInf_sub_shift_bounds_of_gap_pos, Incognita.BatchedBandits.div_two_le_and_one_add_div_le}
\uses{def:constC, def:model, def:kinf}
For $\nu\in\cP$ and $a$ sub-optimal, with $\Delta=\Kinf(\nu_a,\mathbf m^\star)$, $c^\star=\frac{2-\mathbf m^\star}{1-\mathbf m^\star}$ and
$D'_\varepsilon=\Kinf(\nu_a,\mathbf m^\star-2\varepsilon)-\varepsilon$: $\Delta\ge2(\mathbf m^\star-\mathbf m_a)^2>0$, and every
$\varepsilon\in(0,1]$ with $\varepsilon\le\Delta/(4c^\star)$ has
$D'_\varepsilon\ge\Delta-2c^\star\varepsilon\ge\Delta/2\ge\varepsilon^2/8$ and
$(1+\varepsilon)/D'_\varepsilon\le1/\Delta+(2/\Delta+4c^\star/\Delta^2)\varepsilon$.
\end{lemma}

\begin{proof}
\uses{flem:kinf-shift, def:constC, def:kinf}
The paper's proof of Lemma~\ref{lem:constC}(b), with Lemma~\ref{flem:kinf-shift} at $\mu=\mathbf m^\star$ and the shift
$2\varepsilon$; the lower bound on $\Delta$ is the Pinsker bound $\mathrm{kl}(p,\mu)\ge2(\mu-p)^2$.
\end{proof}

\begin{lemma}[the estimates of $\Kinf(\nu_a,\mathbf m^\star)$ when $x>1$]\label{flem:constA-d}
\lean{Incognita.BatchedBandits.dInf_shift_bounds_of_gap_pos}
\uses{def:constA, def:model, def:kinf}
For $\nu\in\cP$, $a$ sub-optimal, $\Delta=\Kinf(\nu_a,\mathbf m^\star)$, $c^\star=\frac{2-\mathbf m^\star}{1-\mathbf m^\star}$, a margin
$\omega\in(0,1]$ and $\varepsilon>0$ with $2\varepsilon\le1-\mathbf m^\star$ and $\varepsilon\le\omega\Delta/(16c^\star)$:
$D_\varepsilon\ge\Delta-2c^\star\varepsilon\ge\frac78\Delta$, $1/D_\varepsilon\le1/\Delta+3c^\star\varepsilon/\Delta^2$,
$K^+_\varepsilon\le\Delta+2c^\star\varepsilon$ and $(1+\omega)D_\varepsilon\ge K^+_\varepsilon$.
\end{lemma}

\begin{proof}
\uses{flem:kinf-shift, def:constA, def:kinf}
The paper's proof of Lemma~\ref{lem:constA}(d), with Lemma~\ref{flem:kinf-shift} at $\mathbf m^\star$ (shift $2\varepsilon$) and at
$\mathbf m^\star+\varepsilon$ (shift $\varepsilon$).
\end{proof}

\begin{lemma}[what ``$\ell$ large'' guarantees when $x=1$]\label{flem:constC-a}
\lean{Incognita.BatchedBandits.mul_div_two_mul_sq_bounds, Incognita.BatchedBandits.log_le_sqrt, Incognita.BatchedBandits.div_six_mul_bounds, Incognita.BatchedBandits.le_inv_pow_max_ceil_log_div_log}
\uses{def:constC}
Let $\delta\in(0,\frac12]$, $\eta>0$, $\varepsilon\in(0,1)$, $\kappa>0$, $w_0\ge0$, an integer $M$ with
$(1-\delta)^{-M}\ge(1+2\eta)\kappa/(2\varepsilon^2)$, $1+\eta\le(1+2\eta)(1-2\delta)$, and $y$ with $y\ge4/\delta^2$,
$y\ge\frac2\delta\log\frac{1+\eta}{2\varepsilon^2}$, $y\ge w_0\kappa/\delta$, $y\ge2w_0$. With $u=(1+\eta)y/(2\varepsilon^2)$:
$1\le u\le e^{(1-\delta)y}$, $(1-\delta)y-\log u\ge(1-2\delta)y>0$, $(1-\delta)^{-M}(1-2\delta)y/\kappa\ge u$,
$w_0\kappa\le\delta y$ and $w_0\le\varepsilon u$. For $\eta>0$ and $\delta=\eta/(6(1+2\eta))$: $0<\delta<\frac1{12}$,
$(1-\delta)(1-2\delta)\ge1-3\delta$ and $(1+2\eta)(1-4\delta)\ge1+\eta$; and
$M=\max(1,\lceil\log X/\log\frac1{1-\delta}\rceil)$ has $(1-\delta)^{-M}\ge X$.
\end{lemma}

\begin{proof}
\uses{def:constC}
The paper's proof of Lemma~\ref{lem:constC}(a), (c) (exponent) and (d): $\log y\le\sqrt y\le\frac\delta2y$ for
$y\ge4/\delta^2$, so $\log u\le\delta y$.
\end{proof}

\begin{lemma}[the series of the consistency test]\label{flem:amb-series}
\lean{Incognita.BatchedBandits.Grid.exists_sum_activated_le, Incognita.BatchedBandits.eventually_rpow_mul_exp_neg_le}
\uses{def:hy, def:grid, def:hy-vocab}
For $c,a>0$ and real $A,b$, eventually $y^Ae^{-cy^a/(\log y)^b}\le y^{-2}$. With the default $\vartheta$, if
$f\ge0$ is bounded on $[1,\infty)$ and eventually at most $Cy^{-2}$, then the sums
$\sum_{\ell\in S}f(\mathsf L_\ell)$ over the finite sets $S$ of activated batches with $t_\ell\ge e^e$ are bounded
uniformly in $S$.
\end{lemma}

\begin{proof}
\uses{lem:sparse, def:hy, def:grid}
$(\log y)^{b+1}=o(y^a)$. For the sums, Lemma~\ref{lem:sparse} bounds them by $\sum_k(k+1)M_k$ with $M_k$ a bound
of $f$ on $[2^k,2^{k+1})$: a constant on the first blocks and $C4^{-k}$ afterwards, and
$(k+1)4^{-k}\le2^{-k}$. This is Step~7 of the paper's proof of Lemma~\ref{lem:amb}, for every function
of the shape of its $f$.
\end{proof}

\begin{lemma}[the bounds $p_A$ and $\bar p_K$ of the consistency test]\label{flem:amb-bounds}
\lean{Incognita.BatchedBandits.mul_exp_neg_div_one_sub_exp_neg_le, Incognita.BatchedBandits.one_add_ceil_mul_exp_neg_div_le, Incognita.BatchedBandits.min_div_two_le_one_sub_exp_neg, Incognita.BatchedBandits.min_div_mul_le_one_sub_exp_neg, Incognita.BatchedBandits.sub_log_le_mul_log, Incognita.BatchedBandits.one_add_ceil_le_mul_sq, Incognita.BatchedBandits.measure_exists_mem_le_pullCount_abs_sub_gt_le_mul_exp_neg_rpow, Incognita.BatchedBandits.measure_exists_mem_le_pullCount_pseudoKinf_lt_or_lt_le_mul_exp_neg, Incognita.BatchedBandits.Grid.eventually_measure_exists_mem_le_pullCount_abs_sub_gt_le, Incognita.BatchedBandits.Grid.eventually_measure_exists_mem_le_pullCount_pseudoKinf_lt_or_lt_le, Incognita.BatchedBandits.mul_exp_neg_le_mul_exp_neg_of_le}
\uses{flem:controlC-run, flem:probasA-primal, flem:ranksC, def:amb, def:kinf, def:model, def:learner, def:grid, def:hy}
(i) For $K\ge0$, $c_\zeta>0$, $y>0$, $\zeta=c_\zeta y^{-1/4}$ and $U\ge(y/4)^\beta$:
$K\cdot2e^{-2\zeta^2(U/2)}/(1-e^{-2\zeta^2})\le p_A(y):=2K(1+\frac{\sqrt y}{2c_\zeta^2})\exp(-c_\zeta^24^{-\beta}y^{\beta-1/2})$.
(ii) For $\Delta>0$, $\mu\in[0,1)$, $y\ge16$, $c_2>0$, $u\ge c_2y$, $\zeta'=\frac\Delta{32}y^{-1/4}$,
$R'=\kappa(\mu)-\log(1-e^{-\zeta'/4})$, $\rho'=\zeta'^2/(2R'^2)$ and $n^{\mathrm g}=\lceil8/(\zeta'(1-e^{-\zeta'/4})(1-\mu))\rceil+1$:
$(1+n^{\mathrm g})e^{-\rho'u}/(1-e^{-\rho'})\le\bar p_K(y):=C_2(1+\frac1{c_1})\,y\log^2y\,\exp(-\frac{c_1c_2\sqrt y}{\log^2y})$, with
$a_0=\min(\frac\Delta{128},1)$, $A_1=\frac14+\frac{\kappa(\mu)+\log(2/a_0)}{\log16}$, $c_1=\frac{\Delta^2}{2048A_1^2}$, $C_2=3+\frac{512}{\Delta a_0(1-\mu)}$.
(iii) Hence, for every learner on $\nu\in\cP$: the trace of $\mathsf{Acc}_\ell^c$ at a set $T$ of cutoffs (some arm $c$ and
$s\in T$ with $N_c(s)\ge U/2$ and $|\hat{\mathbf m}_c(s)-\mathbf m_c|>\zeta$) has probability at most $p_A(y)$; and
the trace of $\mathsf{KC}_\ell^c$ at a finite set $T$ (at the threshold $u$ and the accuracy $\zeta'$) has
probability at most $|T|\,\bar p_K(y)$.
(iv) For $P_0,P_1\ge0$ and $q>0$, $\pi(y)=(P_0+P_1(1+(P_2y)^4))e^{-qy}$ is nonincreasing on $[4/q,\infty)$ (condition (f) of
Definition~\ref{def:amb}, with $q=1+\eta$).
\end{lemma}

\begin{proof}
\uses{flem:controlC-run, flem:probasA-primal, flem:ranksC, def:amb}
Steps~5 and~6 of the paper's proof of Lemma~\ref{lem:amb}: $\frac1{1-e^{-v}}\le1+\frac1v$; $1-e^{-v}\ge\frac{\min(v,1)}2$, so
$1-e^{-\zeta'/4}\ge\frac{a_0}2y^{-1/4}$; then $R'\le A_1\log y$, $\rho'\ge c_1y^{-1/2}/\log^2y$, $1+n^{\mathrm g}\le C_2\sqrt y$ and
$\rho'u\ge c_1c_2\sqrt y/\log^2y$. (iii) is (i) and (ii) applied to the bounds of Lemmas~\ref{flem:controlC-run}(a) at
$u=U/2$ and~\ref{flem:probasA-primal}. With $y=\mathsf L_\ell$, $U=U_{\ell-1}$, $u=n_\ell$, $\mu=\mathbf m^\star$, $\Delta=\Kinf(\nu_a,\mathbf m^\star)$ and
$T=\{t_{\ell-1},t_\ell\}$ these are the paper's bounds on $\PP(\mathsf{Acc}_\ell^c)$ and $\PP(\mathsf{KC}_\ell^c)$, the latter with a
factor $2$. Along a grid with $x_\ell\to1$ and the default $\gamma$, Lemma~\ref{flem:ranksC} gives, for every large $\ell$,
$\gamma_\ell=\mathsf L_\ell^{-1/4}$ and $U_{\ell-1}\ge(\mathsf L_\ell/4)^\beta$: so (iii) holds at $y=\mathsf L_\ell$, $\zeta=c_\zeta\gamma_\ell$ and
$\zeta'=\gamma_\ell\Delta/32$, for every count threshold $u_\ell\ge c_2\mathsf L_\ell$. (iv) For $4/q\le y\le y'$ and $t=y'/y$,
the prefactor grows at most by $t^4$, and $4\log t\le4(t-1)\le qy(t-1)=q(y'-y)$.
\end{proof}

\begin{lemma}[the series of the consistency test, for the paper's $f$]\label{flem:amb-series-f}
\lean{Incognita.BatchedBandits.Grid.exists_sum_activated_rpow_add_mul_exp_neg_le, Incognita.BatchedBandits.rpow_add_mul_sum_exp_neg_le, Incognita.BatchedBandits.exists_rpow_mul_exp_neg_le}
\uses{flem:amb-series, flem:amb-bounds, def:hy, def:grid, def:amb}
With the default $\vartheta$, let $\beta>\frac12$, $k_0,k_1,P_0,P_1,K,D\ge0$, $q,c,c_3>0$, and
$f(y)=\mathfrak c(y)\big(2\pi(y)+p_A(y)+\bar p_K(y)\big)$ for $y\ge16$, $f(y)=0$ below, where $\mathfrak c(y)=y^\beta+k_0+k_1y$,
$\pi(y)=(P_0+P_1(1+(P_2y)^4))e^{-qy}$, $p_A(y)=2K(1+\frac{\sqrt y}{2c^2})\exp(-c^24^{-\beta}y^{\beta-1/2})$ and
$\bar p_K(y)=Dy\log^2y\exp(-c_3\sqrt y/\log^2y)$. Then the sums $\sum_{\ell\in S}f(\mathsf L_\ell)$ over the finite sets $S$ of
activated batches with $t_\ell\ge e^e$ are bounded uniformly in $S$.
\end{lemma}

\begin{proof}
\uses{flem:amb-series}
For $y\ge1$ and $m=\max(\beta,1)$, $f(y)$ is at most a combination with nonnegative coefficients of
$y^{m+4}e^{-qy}$, $y^{m+1}\exp(-c^24^{-\beta}y^{\beta-1/2})$ and $y^{m+2}\exp(-c_3\sqrt y/\log^2y)$ ($\sqrt y\le y$, $\log^2y\le y$).
Each is bounded on $[1,\infty)$ (the exponential is at most $1$) and at most $y^{-2}$ beyond some $y_0$ (Lemma~\ref{flem:amb-series}), so
Lemma~\ref{flem:amb-series} applies to $f$. With $k_0=1+K$, $k_1=16\bar y/\varepsilon^2$, $P_0=2K/(1-e^{-2\varepsilon^2})$,
$P_1=2C_\eta$, $P_2=(1+\eta)/(2\varepsilon^2)$, $q=1+\eta$, $c=c_\zeta$, $D=2C_2(1+1/c_1)$ and $c_3=c_1c_2$, this is Step~7 of the
paper's proof of Lemma~\ref{lem:amb}: once $c_\ell\,\PP(E^{\mathrm{amb}}_\ell(\varepsilon))\le f(\mathsf L_\ell)$, $C_4(\varepsilon)<\infty$.
\end{proof}

\begin{lemma}[the primal and dual formulas of $\Kinf$]\label{flem:primal}
\lean{Incognita.BatchedBandits.dInf_unitLaws_eq_iSup_dual, Incognita.BatchedBandits.exists_dInf_sub_lt_dual, Incognita.BatchedBandits.ofReal_dual_le_dInf, Incognita.BatchedBandits.exists_isMaxOn_dual, Incognita.BatchedBandits.dual_variation, Incognita.BatchedBandits.exists_unitLaws_integral_id_eq_klDiv_le, Incognita.BatchedBandits.sub_two_mul_sq_le_log_one_add, Incognita.BatchedBandits.integrable_of_abs_le_of_mem_unitLaws, Incognita.BatchedBandits.one_sub_mul_sub_bounds, Incognita.BatchedBandits.abs_log_one_sub_mul_sub_le, Incognita.BatchedBandits.integrable_log_one_sub_mul_sub, Incognita.BatchedBandits.integrable_sub_div_one_sub_mul_sub, Incognita.ae_llr_le_neg_log_of_withDensity_le}
\uses{def:kinf}
Let $P$ be a law on $[0,1]$, $\mu\in[0,1)$, $\varphi(\lambda)=\E_P[\log(1-\lambda(X-\mu))]$ and
$D_\lambda=\E_P[(X-\mu)/(1-\lambda(X-\mu))]$ for $\lambda\in[0,\frac1{1-\mu})$. Then:
(i) for $|h|\le c_0/2$, $c_0=1-\lambda(1-\mu)$, $-hD_\lambda-2(h/c_0)^2\le\varphi(\lambda+h)-\varphi(\lambda)\le-hD_\lambda$;
(ii) if $D_\lambda\le0$, the law $Q=g_\lambda P+(-\lambda D_\lambda)\delta_1$, $g_\lambda(x)=1/(1-\lambda(x-\mu))$, is a law on
$[0,1]$ with mean $\mu+D_\lambda(1-\lambda(1-\mu))$ and $\KL(P,Q)\le\varphi(\lambda)$;
(iii) every $\varepsilon>0$ has some $\lambda\in[0,\frac1{1-\mu})$ with $\varphi(\lambda)>\Kinf(P,\mu)-\varepsilon$;
(iv) $\varphi(\lambda)\le\Kinf(P,\mu)$ for every $\lambda\ge0$ with $\lambda(1-\mu)<1$ (every level $\mu$);
(v) hence $\Kinf(P,\mu)=\sup_{\lambda\in[0,\frac1{1-\mu})}\varphi(\lambda)$.
\end{lemma}

\begin{proof}
\uses{flem:kinf-shift, def:kinf}
(i) $1-(\lambda+h)(x-\mu)=(1-\lambda(x-\mu))(1+u)$ with $u=-h(x-\mu)g_\lambda(x)$, $|u|\le1/2$, and
$u-2u^2\le\log(1+u)\le u$. (ii) The law $Q$ is the minimiser of the proof of [HT10, Theorem~8] (the tilt of $P$
by $g_\lambda$, with an atom at $1$): $\E_P[g_\lambda]=1+\lambda D_\lambda$, so $Q$ is a probability; $dP/dQ\le1/g_\lambda$. (iii) If
$\mathbf m(P)>\mu$, $\Kinf=0=\varphi(0)$. Otherwise take $b$ with $1-b(1-\mu)=s$ small and a maximiser $\lambda$ of the
continuous $\varphi$ on $[0,b]$: by (i), $D_\lambda\le0$, and $D_\lambda=0$ if $\lambda<b$; with $D_\lambda\ge-\mu$, the mean of $Q$
of (ii) is at least $\mu-s$, and the mixture $(1-v)Q+v\delta_1$, $v=2s/(1-\mu)$, has mean $>\mu$ and
$\KL\le\varphi(\lambda)-\log(1-v)\le\varphi(\lambda)+2v$ (Lemma~\ref{flem:kinf-shift}). This follows the construction
of the proof of [HT10, Theorem~8], with the first-order conditions taken from (i); it is the direction of
Lemma~\ref{lem:primal} that the upper bound needs: an $\varepsilon$-optimal $\lambda$ of the dual formula, against the
primal formula of the frozen statements. (iv) The variational formula
$\KL(P,Q)\ge\E_P[f]-\E_Q[e^f-1]$ with $f=\log(1-\lambda(x-\mu))$ on $[0,1]$, for which
$\E_Q[e^f-1]=-\lambda(\mathbf m(Q)-\mu)\le0$ when $\mathbf m(Q)>\mu$. (v) From (iii), (iv) and
$\Kinf(P,\mu)<\infty$ (Lemma~\ref{flem:kinf-shift}).
\end{proof}

\begin{lemma}[the deviations of the empirical $\Kinf$ at the primal levels]\label{flem:probasA-primal}
\lean{Incognita.BatchedBandits.measure_le_pullCount_pseudoKinf_lt_dInf_sub_le, Incognita.BatchedBandits.measure_le_pullCount_dInf_add_lt_pseudoKinf_le, Incognita.BatchedBandits.measure_exists_mem_le_pullCount_pseudoKinf_lt_or_lt_le}
\uses{def:evA, def:amb, def:stacks, def:model, def:kinf, def:learner}
For every learner on $\nu\in\cP$, a cutoff $t$, an arm $a$, $\mu\in[0,1)$, $\varepsilon'>0$ and $u>0$: the probability that
$N_a(t)\ge u$ and $\Kinf(\hat\nu_a(t),\mu)<\Kinf(\nu_a,\mu)-\varepsilon'$ is at most $e^{-\rho'u}/(1-e^{-\rho'})$, and
the probability that $N_a(t)\ge u$ and $\Kinf(\hat\nu_a(t),\mu)>\Kinf(\nu_a,\mu)+\varepsilon'$ is at most
$n_ge^{-\rho'u}/(1-e^{-\rho'})$, with the constants of Lemmas~\ref{lem:lowerdev} and~\ref{lem:upperdev}. For a finite set $T$ of
cutoffs, the trace of $\mathsf{KC}^c$ (some $s\in T$ with $N_a(s)\ge u$ and $|\Kinf(\hat\nu_a(s),\mu)-\Kinf(\nu_a,\mu)|>\varepsilon'$) has
probability at most $|T|$ times the sum of the two bounds.
\end{lemma}

\begin{proof}
\uses{flem:primal, flem:kinf-shift, lem:lowerdev, lem:upperdev, lem:probasA}
The bounds of Lemma~\ref{lem:probasA}(c) for its traces hold at a dual level $A$ of one $\lambda$ and at any upper
bound $K^+$ of the dual formula; Lemma~\ref{flem:primal} gives an $\varepsilon'/4$-optimal $\lambda$ with
$A-3\varepsilon'/4>\Kinf(\nu_a,\mu)-\varepsilon'$, and the bound $K^+=\Kinf(\nu_a,\mu)$, a real
number since $\Kinf(\nu_a,\mu)<\infty$ for $\mu<1$ (Lemma~\ref{flem:kinf-shift}).
\end{proof}

\begin{lemma}[the tilts of the sub-optimal arms]\label{flem:tilts-all}
\lean{Incognita.BatchedBandits.exists_tilt_moment_le_forall}
\uses{def:constA, def:model}
For $\nu\in\cP$ there are tilts $\theta_c>0$ of the sub-optimal arms, $c_\theta>0$ and $\varepsilon_\theta>0$ such that every
sub-optimal $c$, $\varepsilon\le\varepsilon_\theta$, $r\le\mathbf m^\star+\varepsilon$ and $\eta\ge0$ with $2\eta\kappa(r)\le c_\theta$ satisfy the
moment conditions of Lemma~\ref{lem:pair} for $\nu_c$ at the level $\mathbf m^\star-3\varepsilon$.
\end{lemma}

\begin{proof}
\uses{lem:theta, lem:constA, def:constA}
Lemma~\ref{lem:theta} for each sub-optimal arm (Lemma~\ref{lem:constA}(c)), and the minima over the finitely
many arms of Definition~\ref{def:constA}.
\end{proof}

\begin{lemma}[the conditions of the rank $\ell_4$]\label{flem:ranksC}
\lean{Incognita.BatchedBandits.Grid.eventually_netSize_bounds, Incognita.BatchedBandits.Grid.eventually_mul_t_le_imedEnd_sub, Incognita.BatchedBandits.Grid.eventually_mul_logEnd_le_netSize, Incognita.BatchedBandits.Grid.eventually_log_le_hyGamma_mul, Incognita.BatchedBandits.Grid.eventually_le_hyGamma_mul_logEnd, Incognita.BatchedBandits.Grid.eventually_le_hyGamma_mul_netSize, Incognita.BatchedBandits.Grid.eventually_logEnd_le_four_mul, Incognita.BatchedBandits.Grid.eventually_div_four_rpow_le_netSize, Incognita.BatchedBandits.Grid.eventually_hyGamma_eq, Incognita.BatchedBandits.Grid.netSize_succ_le, Incognita.BatchedBandits.Grid.logEnd_monotone}
\uses{def:amb, def:water, def:grid, def:hy, def:ass-x}
On a grid with $x_\ell\to1$, for $\beta>1$, the default $\gamma$ and constants $c,c'>0$, $C,C',K'$, $\delta<1$, every
large $\ell$ has: $\gamma_\ell=\mathsf L_\ell^{-1/4}$; $C\mathsf L_\ell\le U_\ell$; $C\log(c\mathsf L_\ell)\le\gamma_\ell c'c\mathsf L_\ell$ and
$e\le c\mathsf L_\ell$; $C\le\gamma_\ell c\mathsf L_\ell$ and $C\le\gamma_\ell cU_\ell$; $\mathsf L_\ell\le4\mathsf L_{\ell-2}$, hence $U_{\ell-1}\ge(\mathsf L_\ell/4)^\beta$; and
$U_{\ell-1}\le2e^{(1-\delta)\mathsf L_\ell}$, $C\mathsf L_\ell\le U_{\ell-1}$, $C'\le U_{\ell-1}$, $K'U_{\ell-1}\le t_{\ell-1}$. These are the
conditions (a)--(e) of Definition~\ref{def:amb} (with $n_\ell=c\mathsf L_\ell$). With the default $\vartheta$, every large $\ell$
also has $s_\ell-t_\ell\ge Kt_\ell$ (the third condition defining $\ell_W$ in Definition~\ref{def:water}).
\end{lemma}

\begin{proof}
\uses{def:amb, def:grid, def:hy}
$U_\ell=\lceil\mathsf L_{\ell-1}^\beta\rceil$ is a power of $\mathsf L$, $t_\ell=e^{\mathsf L_\ell}$, $\log=o(\mathsf L^{3/4})$, and
$\mathsf L_\ell=x_\ell x_{\ell-1}\mathsf L_{\ell-2}$ with $x\le2$ eventually; $t_\ell^{\vartheta_\ell}=e^{\mathsf L_\ell/\log\mathsf L_\ell}\to\infty$.
\end{proof}

\begin{lemma}[the upper bound gives Theorem~\ref{thm:upper}]\label{flem:upper-assembly}
\lean{Incognita.BatchedBandits.isBigO_log_t_of_eventually_div_lt, Incognita.BatchedBandits.isConsistentAtBatchEnds_of_isBigO, Incognita.BatchedBandits.isConsistentAtBatchEnds_of_upper, Incognita.BatchedBandits.tendsto_div_log_t_of_upper, Incognita.BatchedBandits.tendsto_regret_div_log_t_of_tendsto, Incognita.BatchedBandits.hy_isConsistent_of_upper, Incognita.BatchedBandits.hy_tendsto_count_of_upper, Incognita.BatchedBandits.hy_tendsto_regret_of_upper, Incognita.BatchedBandits.eventually_div_lt_of_forall_le_add}
\uses{thm:lower, def:ass-x, def:model, def:learner, def:grid, def:kinf, def:hy}
Let the grid be convergent with limit $x\ge1$ and let a batched learner satisfy
$\limsup_\ell\E_\nu[N_c(t_\ell)]/\log t_\ell\le x^2/\Kinf(\nu_c,\mathbf m^\star)$ for every $\nu\in\cP$
and every sub-optimal arm $c$. Then it is consistent on $\cP$, and for every $\nu\in\cP$ and every
sub-optimal arm $a$, $\E_\nu[N_a(t_\ell)]/\log t_\ell\to x^2/\Kinf(\nu_a,\mathbf m^\star)$ and
$\mathfrak R(t_\ell)/\log t_\ell\to C_x$. In particular, for HY (which is batched on every grid,
Lemma~\ref{flem:hy-congr}), this $\limsup$ bound gives the three claims of
Theorem~\ref{thm:upper}. A sequence with $u_\ell\le A\mathsf L_\ell+C\varepsilon\mathsf L_\ell+R(\varepsilon)$ eventually, for every
$\varepsilon\in(0,\varepsilon_0]$ (the form of Proposition~\ref{prop:fixedC}), has $\limsup_\ell u_\ell/\mathsf L_\ell\le A$.
\end{lemma}

\begin{proof}
\uses{thm:lower, flem:hy-congr, def:learner, def:grid, def:kinf}
The bound gives $\E_\nu[N_c(t_\ell)]=O(\log t_\ell)$ for every $\nu\in\cP$ and sub-optimal $c$, and
$\log t_\ell=o(t_\ell^p)$ for every $p>0$: the learner is consistent. Theorem~\ref{thm:lower} with
logarithmic counts then gives $\liminf_\ell\E_\nu[N_a(t_\ell)]/\log t_\ell\ge
x^2/\Kinf(\nu_a,\mathbf m^\star)$, hence the limit; the regret is the finite sum of the limits
weighted by the gaps (the arms of gap $0$ contribute $0$). This is the end of the paper's proof of
Theorem~\ref{thm:upper} (Subsection~\ref{sec:app-proofupper}), with the upper bound of
Propositions~\ref{prop:dyA} and~\ref{prop:fixedC} as a hypothesis.
\end{proof}

\section{New decompositions of the paper's proofs}

\begin{proof}[Formal route to Lemma~\ref{lem:onebatch}]
\proves{lem:onebatch}
\uses{flem:onebatch-finite, flem:mean-lt-one, def:kinf, def:model, def:grid, def:learner}
Lean: \texttt{Incognita.BatchedBandits.eventually\_lt\_integral\_pullCount\_div\_log\_t\_succ}, and
\texttt{Incognita.P000.one\_batch\_ahead} for the frozen statement. Let
$0\le c<c'<1/\Kinf(\nu_a,\mathbf m^\star)$ (for $c<0$ there is nothing to prove). By the primal
formula there is a law $Q$ on $[0,1]$ with mean $>\mathbf m^\star$ and $c'\KL(\nu_a,Q)<1$; it has
mean $<1$ by Lemma~\ref{flem:mean-lt-one}, so $\nu'=\nu^Q\in\cP$ and every arm $b\ne a$ is
sub-optimal for $\nu'$. This replaces Steps~1 and~3 of the paper (Lemma~\ref{lem:kinf}(v), the
shifted level $\mu'$, and data processing for the binary relative entropy) by the strict
inequality in the primal formula and Bretagnolle--Huber. Let $D=\KL(\nu_a,Q)$, $p=(1-c'D)/2>0$,
$\delta=1/(8K)$. By consistency under $\nu$ and $\nu'$, eventually
$\E_\nu[N_a(t_{\ell+1})]\le\delta t_{\ell+1}^p$ and $\E_{\nu'}[N_b(t_{\ell+1})]\le\delta t_{\ell+1}^p$
for $b\ne a$. Lemma~\ref{flem:onebatch-finite} with $r=t_\ell$, $M=t_{\ell+1}$ (a batched learner
sees the rewards before $t_\ell$ at the rounds $<t_{\ell+1}$) and
$M-\E_{\nu'}[N_a(M)]=\sum_{b\ne a}\E_{\nu'}[N_b(M)]$ give $e^{-\E_\nu[N_a(t_\ell)]D}\le M^{p-1}$,
so $\E_\nu[N_a(t_\ell)]D\ge(1-p)\log t_{\ell+1}$. Since $1-p>c'D\ge0$, this forces $D>0$ and
$\E_\nu[N_a(t_\ell)]>c'\log t_{\ell+1}>c\log t_{\ell+1}$.
\end{proof}

\begin{proof}[Formal route to Theorem~\ref{thm:lower}(i)]
\proves{thm:lower}
\uses{lem:onebatch, def:grid}
Lean: \texttt{Incognita.BatchedBandits.eventually\_lt\_integral\_pullCount\_div\_log\_t} and
\texttt{Incognita.P000.lower\_liminf\_of\_tendsto\_ratio\_one}. As the paper's remark ``every grid'':
$\log t_\ell\le\log t_{\ell+1}$, so $\E[N_a(t_\ell)]/\log t_\ell\ge\E[N_a(t_\ell)]/\log t_{\ell+1}$ for
$\ell\ge1$; the hypothesis $x_\ell\to1$ is not used.
\end{proof}

\begin{proof}[Formal route to Lemma~\ref{lem:twobatches}(ii)]
\proves{lem:twobatches}
\uses{flem:twobatches, def:grid}
Lemma~\ref{flem:twobatches} is part (ii) at the index $\ell+1$. Part (i), the form with tightness
in probability, and Lemmas~\ref{lem:tilts}, \ref{lem:com} and~\ref{lem:typ} are not formalised:
the route of Lemma~\ref{flem:twobatches} does not need them.
\end{proof}

\begin{proof}[Formal route to Theorem~\ref{thm:lower}(ii)--(iii) and Theorem~\ref{thm:main}(a)]
\proves{thm:lower}
\uses{flem:twobatches, flem:caseA, flem:caseB, flem:liminf-sum, flem:regret-one,
  flem:regret-caseB, lem:onebatch, def:grid, def:ass-x, def:model, def:learner, def:kinf}
Lean: \texttt{Incognita.BatchedBandits.eventually\_lt\_integral\_pullCount\_div\_log\_t\_of\_isBigO},
\texttt{Incognita.BatchedBandits.eventually\_lt\_regret\_div\_log\_t\_of\_forall}, and the frozen
statements \texttt{Incognita.P000.lower\_liminf}, \texttt{lower\_limsup},
\texttt{lower\_regret\_liminf}, \texttt{lower\_regret\_limsup}, for every $x\ge1$. Counts, $\liminf$:
for $x=1$, Theorem~\ref{thm:lower}(i) on every grid; if $\E[N_a(t_\ell)]=O(\log t_\ell)$, then
$\E[N_a(t_{\ell+1})]=O(\log t_{\ell+2})$, Lemma~\ref{flem:twobatches} gives
$\liminf\E[N_a(t_\ell)]/\log t_{\ell+2}\ge1/\Kinf$ and Lemma~\ref{flem:caseA} the factor $x^2$
(Case A of the paper). Counts, $\limsup$: if the counts are logarithmic, the $\liminf$ claim;
otherwise Lemma~\ref{flem:caseB} (Case B). Regret, $\liminf$: for $x=1$
Lemma~\ref{flem:regret-one}; otherwise the bound for the counts of each sub-optimal arm and
Lemma~\ref{flem:liminf-sum}. Regret, $\limsup$: if every sub-optimal arm has logarithmic counts,
the $\liminf$ claim; otherwise Lemma~\ref{flem:regret-caseB}.
\end{proof}

\begin{proof}[Formal route to Theorem~\ref{thm:main}(b), HY is a learner of $\mathcal M_0$]
\proves{thm:main}
\uses{flem:hy-congr, flem:hy-grid, def:M0, def:hy, def:grid}
Lean: \texttt{Incognita.BatchedBandits.infoModel\_zero\_hy} and
\texttt{Incognita.P000.hy\_infoModel\_zero}. Condition (a) of Definition~\ref{def:M0} is
Lemma~\ref{flem:hy-congr}; condition (b) with $k=0$ is Lemma~\ref{flem:hy-grid} ($t_0=0$ on every
grid). The hypotheses $w_0\ge0$, $\beta\ge2$ are not used.
\end{proof}

\end{document}
